\subsection{希尔伯特空间的外希尔伯特和}

命题 1. --- 设 \((\mathrm{E}_i)_{i\in \mathbf{I}}\) 为一族希尔伯特空间，\(\mathbf{P}\) 为乘积向量空间 \(\prod_{i\in \mathbf{I}}\mathrm{E}_i\)，\(\mathbf{E}\) 为 \(\mathbf{P}\) 中由所有家庭 \(\mathbf{x} = (x_i)_{i\in \mathbf{I}}\)（\(\sum_{i\in \mathbf{I}}\| x_i\|^2\) 为有限）组成的子集。

\begin{enumerate}
\def\labelenumi{\alph{enumi})}
\item
  E 是 \(\mathbf{P}\) 的向量子空间。
\item
  对 E 中的每个 \(\mathbf{x} = (x_i)_{i\in \mathbf{I}}\) 与 \(\mathbf{y} = (y_i)_{i\in \mathbf{I}}\)，族 \((\langle x_i|y_i\rangle)_{i\in \mathbf{I}}\) 可和。若令 \(\langle \mathbf{x}|\mathbf{y}\rangle = \sum_{i\in \mathbf{I}}\langle x_i|y_i\rangle\)，就在 E 上定义了一个正的非退化埃尔米特型。
\item
  对于如此定义的内积，\(\mathbf{E}\) 是希尔伯特空间；\(\mathbf{S}\) 是诸 \(\mathbf{E}_i\) 的直和，它在 \(\mathbf{E}\) 中稠密。
\end{enumerate}

对 E 中的 \(\mathbf{x} = (x_{i})_{i\in \mathbf{I}}\) 与 \(\mathbf{y} = (y_i)_{i\in \mathbf{I}}\)，我们有

\[
\left\| x _ {i} + y _ {i} \right\| ^ {2} \leqslant 2 \left(\left\| x _ {i} \right\| ^ {2} + \left\| y _ {i} \right\| ^ {2}\right),
\]

因而 \(\mathbf{x} + \mathbf{y} = (x_i + y_i)_{i\in I}\) 属于 E。这就证明了 \(a)\)。

由 Cauchy-Schwarz 不等式，我们有

\[
\left| \left\langle x _ {i} \mid y _ {i} \right\rangle \right| \leqslant \| x _ {i} \| \cdot \| y _ {i} \| \leqslant \frac {1}{2} \left(\left\| x _ {i} \right\| ^ {2} + \left\| y _ {i} \right\| ^ {2}\right)
\]

因而 \(\sum_{i\in \mathbf{I}}|\langle x_i|y_i\rangle | < +\infty\)。若 \(x\neq 0\)，我们有 \(\langle \mathbf{x}|\mathbf{x}\rangle = \sum_{i\in \mathbf{I}}\| x_i\|^2 >0\)，于是断言 \(b)\) 得证。

我们回顾，S 是 P 中由所有满足如下条件的家庭 \(\mathbf{x} = (x_{i})_{i \in \mathbf{I}}\) 组成的子空间：\(i \in I\) 中使 \(x_{i} \neq 0\) 的那些所成的集是有限的。由此立即推出 S 在 E 中稠密；于是只需证明 E 对拓扑 \(T_{1}\)（由范数 \(\|x\| = \langle x|x\rangle^{1/2}\) 得到）是完备的。设 \(T_{2}\) 为 \(\prod_{i \in I} E_{i}\) 上的乘积拓扑在 E 上诱导的拓扑。对每个 r \textgreater{} 0，设 \(B_{r}\) 为 \(x \in E\) 中所有满足 \(\|x\| \leqslant r\) 的元素组成的集。此关系式蕴含：对 I 的每个有限子集 J 都有 \(\sum_{i \in J} \|x_{i}\|^{2} \leqslant r^{2}\)，因而 \(B_{r}\) 是 \(\prod_{i \in I} E_{i}\) 的闭子集，从而也是完备的。于是 E 对 \(T_{1}\) 完备这一事实可由 GT, III, §3, No.~5, 命题 10 的推论 2 得出。

定义 1. --- 设 \((\mathrm{E}_{i})_{i\in\mathrm{I}}\) 为一族希尔伯特空间。命题 1 中定义的希尔伯特空间 E 称为族 \((\mathrm{E}_{i})_{i\in\mathrm{I}}\) 的外希尔伯特和，记作 \(\bigoplus_{i\in I} E_{i}\) 或 \(\bigoplus_{i\in I} E_{i}^{1}\)。

\(^{1}\) 应注意不要把这一记号与诸空间 \(E_{i}\) 的 « 代数 » 直和的记号（A, II, § 1, No.~6）相混淆。

设 \(f_{i}\) 为从 \(E_{i}\) 到 \(E\) 中的映射，它把 \(z \in E_{i}\) 变为满足如下条件的元素 \((x_{k}) \in E\)：\(x_{k} = 0\) 对所有 \(k \neq i\) 成立，且 \(x_{i} = z\)；显然 \(f_{i}\) 是从希尔伯特空间 \(E_{i}\) 到 \(E\) 的一个闭向量子空间上的同构。我们称 \(f_{i}\) 为从 \(E_{i}\) 到 \(E\) 中的典范映射，并且通常把 \(E_{i}\) 与它在此同构下的、\(E\) 中的像等同起来。在这一约定下，\(E_{i}\) 与 \(E_{k}\) 在 \(E\) 中正交（当 \(i \neq k\) 时），而 \(E\) 是由诸子空间 \(E_{i}\) 的并生成的闭向量子空间。

当 I 为有限时，E 是诸 \(E_{i}\) 的直和；由于从 E 到 \(E_{i}\) 上的典范投影算子对所有 \(i \in I\) 都连续，E 也是诸 \(E_{i}\) 的拓扑直和（GT, III, § 6, No.~2, 命题 2）。若 \(I = \left(1, n\right)\)，我们也用 \(E_{1} \oplus E_{2} \oplus \ldots \oplus E_{n}\) 来代替 \(\bigoplus_{i=1}^{n} E_{i}\)。

例. --- 设 E 为希尔伯特空间，I 为指标集。用 \(\ell_{\mathrm{E}}^{2}(\mathrm{I})\) 表示族 \((\mathrm{E}_{i})_{i\in\mathrm{I}}\)（其中 \(E_{i}=E\) 对所有 \(i\in I\) 成立）的外希尔伯特和。换言之，\(\ell_{\mathrm{E}}^{2}(\mathrm{I})\) 是由 E 中元素的所有家庭 \(\mathbf{x}=(x_{i})_{i\in\mathrm{I}}\)（满足 \(\sum_{i\in I}\|x_{i}\|^{2}<+\infty\)）组成的空间，赋以内积 \(\langle x|y\rangle=\sum_{i\in I}\langle x_{i}|y_{i}\rangle\)（以 I 为指标集的、E 中元素之平方可和家庭所成的空间）。我们令 \(\ell^{2}(\mathrm{I})=\ell_{\mathrm{K}}^{2}(\mathrm{I})\)。

\subsection{希尔伯特空间的正交子空间的希尔伯特和}

定义 2. --- 称希尔伯特空间 E 是其一族闭向量子空间 \((\mathbf{E}_{i})_{i\in I}\) 的希尔伯特和，如果：

\begin{enumerate}
\def\labelenumi{\arabic{enumi})}
\item
  对 I 中两个不同的指标 \(i\)、\(k\)，子空间 \(\mathrm{E}_i\) 与 \(\mathrm{E}_k\) 在 E 中正交；
\item
  由诸 \(\mathbf{E}_i\) 的并生成的闭向量子空间就是 E。
\end{enumerate}

定理 1. --- 设 E 为希尔伯特空间，它是一族闭向量子空间 \((\mathrm{E}_{i})_{i\in\mathbf{I}}\) 的希尔伯特和。存在唯一的、从 E 到 \(\bigoplus_{i\in I}E_{i}=F\)（族 \((\mathrm{E}_{i})\) 的外希尔伯特和）上的同构 f，使得对所有 \(i\in I\)，\(f\) 在 \(\mathbf{E}\) 上的限制就是典范映射 \(f_{i}\)（从 \(\mathbf{E}_i\) 到 \(\mathbf{F}\) 中）。

设 \(S \subset F\) 为诸 \(E_i\) 的 « 代数 » 直和，并设 \(g\) 为线性映射 \((x_i)_{i \in I} \mapsto \sum_{i \in I} x_i\)（从 \(S\) 到 \(E\) 中）。我们将证明 \(g\) 是从准希尔伯特空间 S 到（准希尔伯特）子空间 \(g(S)\)（E 的由诸 \(E_{i}\) 的并生成的子空间）上的同构：因为，对两个元素 \(\mathbf{x} = (x_{i})_{i \in \mathrm{I}}\)、\(\mathbf{y} = (y_{i})_{i \in \mathrm{I}}\)，我们有

\[
\langle g (\mathbf {x}) | g (\mathbf {y}) \rangle = \left\langle \sum_ {i \in \mathrm{I}} x _ {i} \right| \sum_ {i \in \mathrm{I}} y _ {i} \rangle = \sum_ {(i, k) \in \mathrm{I} \times \mathrm{I}} \left\langle x _ {i} \mid y _ {k} \right\rangle .
\]

但若 \(i \neq k\)，则由假设 \(\langle x_i | y_k \rangle = 0\)，因而

\[
\langle g (\mathbf {x}) | g (\mathbf {y}) \rangle = \sum_ {i \in \mathbf {I}} \langle x _ {i} | y _ {i} \rangle = \langle \mathbf {x} | \mathbf {y} \rangle ;
\]

这就证明了我们的断言。由于 S 在 F 中稠密且 \(g(S)\) 在 E 中稠密，同构 g 延拓为从 F 到 E 上的同构 \(\overline{g}\)（V, p.~8, 推论）。显然 \(\overline{g}\) 的逆同构 f 就是所求的映射；其唯一性由以下事实得出：E 中由诸 \(E_{i}\) 的并生成的闭子空间就是 E 自身。

当 E 是一族子空间 \((\mathrm{E}_{i})_{i\in I}\) 的希尔伯特和时，我们常借助同构 f 把 E 与诸 \(E_{i}\) 的外希尔伯特和 F 等同起来。~若集 I 为有限，则说 E 是族 \((\mathrm{E}_{i})_{i\in I}\) 的希尔伯特和，就意味着诸 \(E_{i}\) 两两正交，且向量空间 E 是子空间族 \((\mathrm{E}_{i})_{i\in I}\) 的直和。

推论 1. --- 设 E 为希尔伯特空间，它是一族闭向量子空间 \((\mathbf{E}_{i})_{i\in\mathbf{I}}\) 的希尔伯特和；对所有 \(i\in I\)，设 \(p_{E_{i}}\) 为从 E 到 \(E_{i}\) 上的正交投影算子（V, p.~13）。

\begin{enumerate}
\def\labelenumi{\alph{enumi})}
\tightlist
\item
  对所有 \(x \in \mathbb{E}\)，族 \((\| p_{\mathbf{E}_i}(x)\|^{2})_{i\in \mathbf{I}}\) 在 \(\mathbf{R}\) 中可和，族 \((p_{\mathbf{E}_i}(x))_{i\in \mathbf{I}}\) 在 \(\mathbb{E}\) 中可和，并且我们有
\end{enumerate}

\[
\left\| x \right\| ^ {2} = \sum_ {i \in \mathbf {I}} \left\| p _ {\mathrm{E} _ {i}} (x) \right\| ^ {2}, \quad x = \sum_ {i \in \mathbf {I}} p _ {\mathrm{E} _ {i}} (x).
\]

\begin{enumerate}
\def\labelenumi{\alph{enumi})}
\setcounter{enumi}{1}
\item
  反之，若 \((x_{i})_{i\in \mathbf{I}}\) 是 \(\mathbf{E}\) 中元素的一族，使得 \(x_{i}\in E_{i}\) 对所有 \(i\in \mathbf{I}\) 成立，且 \(\sum_{i\in \mathbf{I}}\| x_i\|^2 < + \infty\)，则此族可和，且其和 \(x\) 是 \(\mathbf{E}\) 中唯一的满足 \(p_{\mathrm{E}_i}(x) = x_i\)（对所有 \(i\in \mathbf{I}\)）的点。
\item
  对每一对点 \(x, y\)（属于 \(\mathbf{E}\)），我们有
\end{enumerate}

\[
\langle x | y \rangle = \sum_ {i \in \mathbf {I}} \left\langle p _ {\mathrm{E} _ {i}} (x) \mid p _ {\mathrm{E} _ {i}} (y) \right\rangle .
\]

这些性质对诸 \(\mathbf{E}_i\) 的外希尔伯特和而言事实上是显然的，并可经由同构转移到 \(\mathbf{E}\) 上。

推论 2. --- 设 E 为分离的准希尔伯特空间，\((\mathrm{E}_{i})_{i\in\mathrm{I}}\) 为 E 的一族完备向量子空间，使得对 I 中每一对不同的指标 i、k，子空间 \(E_{i}\) 与 \(E_{k}\) 正交。设 V 为 E 中由诸 \(E_{i}\) 的并生成的闭向量子空间。对每个 \(i\in I\)，设 \(p_{E_{i}}\) 为从 E 到 \(E_{i}\) 上的正交投影算子。设 \(x\in E\)。

\begin{enumerate}
\def\labelenumi{\arabic{enumi})}
\item
  我们有 \(\sum\left\|p_{E_{i}}(x)\right\|^{2}\leqslant\left\|x\right\|^{2}\)。
\item
  下列条件等价：a) \(x \in \mathbf{V}\)；b) \(\sum_{i \in I} \| p_{\mathrm{E}_i}(x) \|^2 = \| x \|^2\)；c) 族 \((p_{\mathrm{E}_i}(x))_{i \in I}\) 在 E 中可和，且 \(x = \sum_{i \in I} p_{\mathrm{E}_i}(x)\)。
\item
  假设 V 完备。则族 \((p_{\mathrm{E}_i}(x))_{i\in \mathbf{I}}\) 在 E 中可和，且
\end{enumerate}

\[
p _ {\mathbf {V}} (x) = \sum_ {i \in \mathbf {I}} p _ {\mathrm{E} _ {i}} (x), \quad \left\| p _ {\mathbf {V}} (x) \right\| ^ {2} = \sum_ {i \in \mathbf {I}} \left\| p _ {\mathrm{E} _ {i}} (x) \right\| ^ {2},
\]

其中 \(p_{V}\) 表示从 E 到 V 上的正交投影算子。

设 \(\hat{\mathbf{E}}\) 为 \(\mathbf{E}\) 的希尔伯特空间完备化；我们把 \(\mathbf{E}\) 与 \(\hat{\mathbf{E}}\) 的一个稠密子空间等同起来；诸 \(\mathbf{E}_i\) 因完备而是 \(\hat{\mathbf{E}}\) 的闭子空间。闭包 \(\overline{\mathbf{V}}\)（\(\mathbf{V}\) 在 \(\hat{\mathbf{E}}\) 中的闭包）是 \(\hat{\mathbf{E}}\) 中由诸 \(\mathbf{E}_i\) 的并生成的闭向量子空间，且 \(\mathbf{V} = \overline{\mathbf{V}} \cap \mathbf{E}\)。空间 \(\hat{\mathbf{E}}\) 是诸 \(\mathbf{E}_i\) 与子空间 \(\mathbf{W}\)（即 \(\overline{\mathbf{V}}\) 在 \(\hat{\mathbf{E}}\) 中的正交补）的希尔伯特和；令 \(x_0 = p_{\mathrm{W}}(x)\) 与 \(x_i = p_{\mathrm{E}_i}(x)\)（对所有 \(i \in \mathbf{I}\)）。由推论 1，我们有 \(\| x\|^2 = \| x_0\|^2 + \sum_{i \in \mathbf{I}} \| x_i\|^2\)，且 \(x = x_0 + \sum_{i \in \mathbf{I}} x_i\)（在 \(\hat{\mathbf{E}}\) 中）。由此即得断言 1)，以及 2) 的条件 \(b)\) 与 \(c)\) 等价于条件 \(x_0 = 0\)、从而等价于条件 \(x \in \mathbf{V}\) 这一事实。最后，若 \(\mathbf{V}\) 完备，并令 \(x' = p_{\mathrm{V}}(x)\)，我们有 \(x' - x_i = (x - x_i) - (x - p_{\mathrm{V}}(x))\)，因而 \(x' - x_i\) 与 \(\mathbf{E}_i\) 正交，于是 \(x_i = p_{\mathrm{E}_i}(x')\) 对所有 \(i \in \mathbf{I}\) 成立；此时只需对向量 \(x'\) 应用性质 2) 即可。

注. --- 设 E 为分离的准希尔伯特空间，\((\mathrm{V}_{i})_{i\in\mathrm{I}}\) 为 E 的一族向量子空间，使得对每一对不同的指标 i、k，子空间 \(V_{i}\) 与 \(V_{k}\) 正交。于是，对每个 \(k\in I\)，\(V_{k}\) 与闭向量子空间 \(W_{k}\)（由诸 \(V_{i}\)（\(i\neq k\)）的并生成）的交化为 0；因为，若 x 既属于 \(V_{k}\) 又属于 \(W_{k}\)，则它与所有 \(V_{i}\)（\(i\neq k\)）正交，从而与 \(W_{k}\) 正交。特别地，x 与自身正交，故为零。

命题 2. --- 设 E 为希尔伯特空间，\((\mathrm{V}_{\lambda})_{\lambda\in\mathrm{L}}\) 为 E 的一族闭向量子空间；对每个 \(\lambda\in L\)，设 \((\mathrm{W}_{\lambda\mu})_{\mu\in\mathrm{M}_{\lambda}}\) 为 \(V_{\lambda}\) 的一族闭向量子空间，使得 \(V_{\lambda}\) 是由此族之并生成的闭向量子空间。为使 E 是族 \((\mathrm{W}_{\lambda\mu})_{\lambda\in\mathrm{L},\mu\in\mathrm{M}_{\lambda}}\) 的希尔伯特和，必要且充分的是：E 是族 \((\mathrm{V}_{\lambda})_{\lambda\in\mathrm{L}}\) 的希尔伯特和，且对每个 \(\lambda\in L\)，\(V_{\lambda}\) 是族 \((\mathrm{W}_{\lambda\mu})_{\mu\in\mathrm{M}_{\lambda}}\) 的希尔伯特和（« 希尔伯特和的结合性 »）。为证条件必要，只需看出 \(V_{\alpha}\) 与 \(V_{\beta}\) 正交（若 \(\alpha\neq\beta\)）。但是，\(W_{\alpha\mu}\)（\(\mu\in M_{\alpha}\)）的每个元素与所有 \(W_{\beta\nu}\)（\(\nu\in M_{\beta}\)）正交，从而与它们所生成的闭向量子空间 \(V_{\beta}\) 正交；同样的论证进而表明 \(V_{\beta}\) 的每个元素与 \(V_{\alpha}\) 正交，因为它与所有 \(W_{\alpha\mu}\)（\(\mu\in M_{\alpha}\)）正交。

为证条件充分，只需验证：若条件满足，则 \(\mathrm{E}\) 等于闭向量子空间 \(\mathrm{F}\)（由诸 \(\mathbf{W}_{\lambda \mu}\)（\(\lambda \in \mathbf{L}\)，\(\mu\in M_{\lambda}\)）的并生成）；而对每个 \(\lambda\in L\)，F 包含由诸 \(W_{\lambda\mu}\)（其中 \(\mu\in M_{\lambda}\)）的并生成的闭向量子空间，也就是说 F 包含 \(V_{\lambda}\)；于是 F 是由诸 \(V_{\lambda}\) 的并生成的闭向量子空间，而由假设它就是 E。

\subsection{标准正交族}

定义 3. --- 在准希尔伯特空间中，称向量族 \((e_{i})_{i\in\mathbf{I}}\) 为正交的，如果 \(e_{i}\) 与 \(e_{k}\)（对所有 \(i\neq k\)）正交；称其为标准正交的，如果此外 \(\|e_{i}\|=1\) 对所有 \(i\in I\) 成立。

E 的子集 S，若由 S 到其自身的恒等映射所定义的族是标准正交的，则称之为标准正交集。若 \((e_{i})_{i\in I}\) 是标准正交族，则映射 \(i\mapsto e_{i}\) 是单射；于是我们可以不加区分地谈论标准正交族或标准正交集。

若 \((e_i)_{i\in I}\) 是标准正交族，则完备的一维向量子空间 \(\mathbf{D}_i = \mathbf{K}e_i\) 两两正交。对每个 \(x\in E\)，\(x\) 在 \(\mathbf{D}_i\) 上的正交投影是 \(\lambda_i e_i\)，其中 \(\langle e_i|x - \lambda_i e_i\rangle = 0\)，由此得 \(\langle e_i|x\rangle = \lambda_i\langle e_i|e_i\rangle = \lambda_i\)。将 No.~2 的结果应用于诸子空间 \(\mathbf{D}_i\)，便得到下列命题：

命题 3. --- 在分离的准希尔伯特空间 E 中，每个标准正交族都是拓扑无关的。

我们注意到，这一性质可由拓扑无关族的刻画（IV, p.~1 与 II, p.~43, 推论 2）立即得出，因为按 K 等于 R 或 C，E 的对偶分别与 E 的完备化或 E 的共轭空间等同（V, p.~17, 注）。

命题 4. --- 设 E 为分离的准希尔伯特空间，\((e_{i})_{i\in I}\) 为 E 中的标准正交族，V 为 E 中由诸 \(e_{i}\) 生成的闭向量子空间。

\begin{enumerate}
\def\labelenumi{\arabic{enumi})}
\tightlist
\item
  对每个 \(x \in E\)，我们有
\end{enumerate}

\[
\sum_ {i \in \mathbf {I}} \left| \langle e _ {i} | x \rangle \right| ^ {2} \leqslant \| x \| ^ {2}\tag{1}
\]

（Bessel 不等式）；这里，所有 \(i \in \mathbf{I}\) 中使 \(\langle e_i|x\rangle \neq 0\) 者所成的集是可数的。此外，下列条件等价：a) \(x \in \mathbf{V}\)；b) \(\|x\|^2 = \sum_{i \in \mathbf{I}} |\langle e_i|x\rangle|^2\)；c) 族 \(\langle e_i|x\rangle\) . \(e_i\) 在 E 中可和，且 \(x = \sum_{i \in \mathbf{I}} \langle e_i|x\rangle\) . \(e_i\)。

\begin{enumerate}
\def\labelenumi{\arabic{enumi})}
\setcounter{enumi}{1}
\item
  若 \(\mathbf{V}\) 完备，则由所有 \(\langle e_i|x\rangle\) . \(e_i\) 组成的族在 \(\mathbf{E}\) 中（对所有 \(x\in\mathbf{E}\)）可和，且 \(\sum_{i\in\mathbf{I}}\langle e_i|x\rangle\) . \(e_i=p_{\mathbf{V}}(x)\)，\(\sum_{i\in\mathbf{I}}|\langle e_i|x\rangle|^2=\|p_{\mathbf{V}}(x)\|^2\)。
\item
  假设 V 完备。对每个标量族 \((\lambda_{i})_{i\in\mathbf{I}}\)（满足 \(\sum_{i\in I}|\lambda_{i}|^{2}<+\infty\)），存在唯一的点 \(x\in V\)，使得 \(\langle e_{i}|x\rangle=\lambda_{i}\) 对所有 \(i\in I\) 成立。若 \((\mu_{i})_{i\in I}\) 是第二个满足 \(\sum_{i\in I}\left|\mu_{i}\right|^{2}<+\infty\) 的标量族，且 \(y\in V\) 满足 \(\langle e_{i}|y\rangle=\mu_{i}\)（对所有 \(i\in I\)），则 \(\langle x|y\rangle=\sum_{i\in I}\overline{\lambda}_{i}\mu_{i}\)。
\end{enumerate}

命题 5. --- 设 \((e_i)_{i \in I}\) 为分离的准希尔伯特空间 E 中的标准正交族。下列性质等价：

\begin{enumerate}
\def\labelenumi{\alph{enumi})}
\item
  族 \((e_i)\) 是完全的；
\item
  对每个 \(x \in E\)，族 \(\langle e_{i}|x\rangle\) . \(e_{i}\) 在 E 中可和，且 \(x = \sum_{i \in I} \langle e_{i}|x\rangle\) . \(e_{i}\)；
\item
  对每个 \(x \in E\)，
\end{enumerate}

\[
\left\| x \right\| ^ {2} = \sum_ {i \in \mathbf {I}} \left| \langle e _ {i} | x \rangle \right| ^ {2}\tag{2}
\]

（Parseval 等式）。

当 E 是希尔伯特空间时，这些条件也等价于：

\begin{enumerate}
\def\labelenumi{\alph{enumi})}
\setcounter{enumi}{3}
\tightlist
\item
  关系式 \(\langle e_i|x\rangle = 0\)（对所有 \(i\in \mathrm{I}\)）蕴含 \(x = 0\)。
\end{enumerate}

条件 a)、b)、c) 的等价性由命题 4 立即得出。当 E 是希尔伯特空间时，条件 a) 与 d) 的等价性由 V, p.~16 的推论 1 得出。

定义 4. --- 分离的准希尔伯特空间 E 中的标准正交且完全的族称为 E 的标准正交基。

分离的准希尔伯特空间 E 的标准正交基也是 E 的完备化的标准正交基。

设 \((e_{i})_{i\in I}\) 为 E 的标准正交基；对每个 \(x\in E\)，数 \(\langle e_{i}|x\rangle\) 滥用术语称为 x 关于基 \((e_{i})\) 的坐标。对 E 中每个 x 与 y，我们有

\[
\langle x | y \rangle = \sum_ {i \in \mathbf {I}} \overline {{\langle e _ {i} | x \rangle}} \langle e _ {i} | y \rangle .\tag{3}
\]

E 的标准正交基一般并不是 A, II, p.~25 中所定义的意义下 E 在域 K 上的基；为避免任何混淆，我们把上述引文意义下准希尔伯特空间 E 的基称为 E 在 K 上的代数基。

设 E 与 F 为两个分离的准希尔伯特空间，u 为从 E 到 F 中的连续线性映射。设 \((e_{i})_{i\in I}\)（相应地，\((f_{j})_{j\in J}\)）为 E（相应地，F）的标准正交基。令

\[
u _ {j i} = \langle f _ {j} | u (e _ {i}) \rangle
\]

其中 \(i \in \mathbf{I}\)，\(j \in \mathbf{J}\)。族 \((u_{ji})_{(i,j) \in \mathbf{I} \times \mathbf{J}}\) 称为 \(u\) 关于标准正交基 \((e_i)\) 与 \((f_j)\) 的矩阵。设 \(x \in \mathbf{E}\) 且 \(y = u(x)\)；若分别用 \(\xi_i = \langle e_i | x \rangle\) 与 \(\eta_j = \langle f_j | y \rangle\) 表示 \(x\) 与 \(y\) 的坐标，则 \(\eta_j = \sum_{i \in \mathbf{I}} u_{ji} \xi_i\) 对所有 \(j \in \mathbf{J}\) 成立。当 \((e_i)\) 是 \(E\) 的代数基且 \((f_j)\) 是 \(F\) 的代数基时，我们的定义与 A, II, § 10, No.~4 中的定义一致。

例. --- 设 E 为 R 上复值连续函数的空间，这些函数满足 \(f(x + n) = f(x)\)（\(x \in R\)，\(n \in Z\)）。我们赋予 E 由下式定义的内积

\[
\langle f | g \rangle = \int_ {0} ^ {1} \overline {{{{f (t)}}}} g (t) d t.
\]

于是 E 是分离的准希尔伯特空间，但不完备。对每个整数 \(n \in Z\)，令 \(e_{n}(x) = \mathbf{e}(nx)\)。立即可以看出族 \((e_{n})_{n \in \mathbf{Z}}\) 在 E 中标准正交。此外，E 上的一致收敛拓扑细于由范数 \(\|f\|_{2} = \langle f | f \rangle^{1/2}\) 得到的拓扑。族 \((e_{n})_{n \in \mathbf{Z}}\) 在 E 中对一致收敛是完全的（GT, X, § 4, No.~4），从而在准希尔伯特空间 E 中更是完全的。因此 \((e_{n})_{n \in \mathbf{Z}}\) 是 E 的标准正交基。

\subsection{标准正交化}

定理 2. --- 对希尔伯特空间 E 中的每个标准正交集 L，存在 E 的包含 L 的标准正交基 B。

事实上，设 D 为 E 的所有标准正交子集按包含关系线性排序所成的族；立即可以看出此族具有有限特征（S, III, § 4, No.~5）。于是由 S, III, § 4, No.~5 的定理 1，D 中存在包含 L 的极大族 B。剩下要证 B 是完全集。若不然，将存在一个与 B 的所有向量都正交的向量 \(y \neq 0\)（V, p.~22, 命题 5），用适当的标量乘 y，可设 \(\|y\| = 1\)；那么 \(B \cup \{y\}\) 将是异于 B 且包含 B 的标准正交集；这与 B 的定义矛盾；定理得证。

推论 1. --- 在每个希尔伯特空间中都存在标准正交基。

只需对 L = ∅ 的情形应用定理 2。

推论 2. --- 每个希尔伯特空间都同构于某个空间 \(\ell^2 (\mathbf{I})\)。

更确切地说，设 \((e_{i})_{i\in I}\) 为希尔伯特空间 E 的标准正交基。由命题 4（V, p.~21）与命题 5（V, p.~22），由下式定义的映射 \(\phi\)

\[
\phi (x) = \big (\langle e _ {i} | x \rangle \big) _ {i \in \mathrm{I}}\tag{4}
\]

是从 E 到 \(\ell^2 (\mathrm{I})\) 上的希尔伯特空间同构。逆同构 \(\psi\) 由下式定义

\[
\psi ((\lambda_ {i}) _ {i \in \mathbf {I}}) = \sum_ {i \in \mathbf {I}} \lambda_ {i} e _ {i}.\tag{5}
\]

命题 6. --- 设 E 为分离的准希尔伯特空间，\((a_{n})_{n\in\mathbb{I}}\)（I 是 N 的以 1 为起点的区间）为 E 中向量的一个可数的（有限或无限）无关族。则存在唯一的 E 中的标准正交族 \((e_{n})_{n\in\mathbb{I}}\)，具有下列性质：

\begin{enumerate}
\def\labelenumi{\arabic{enumi})}
\item
  对每个整数 \(p \in \mathbf{I}\)，\(\mathbf{E}\) 的由 \(e_1, e_2, \ldots, e_p\) 生成的向量子空间与 \(\mathbf{E}\) 的由 \(a_1, a_2, \ldots, a_p\) 生成的向量子空间相同；
\item
  对每个整数 \(p \in \mathbf{I}\)，数 \(\langle a_p | e_p \rangle\) 是实数且 \(> 0\)。
\end{enumerate}

事实上，设 \(V_{n}\) 为由 \(a_{1}, a_{2}, ..., a_{n}\) 生成的（n 维）子空间。若 \(n + 1 \in I\) 且 \(b_{n+1} = a_{n+1} - p_{\mathrm{V}_{n}}(a_{n+1})\)（其中 \(p_{V_{n}}\) 是到完备子空间 \(V_{n}\) 上的正交投影算子），则直线 \(Kb_{n+1}\) 是 \(V_{n}\) 在 \(V_{n+1}\) 中的正交（补）。若诸 \(e_{n}\) 满足命题的条件 1)，则必须有 \(e_{n+1} = \lambda b_{n+1}\)；条件 \(\|e_{n+1}\| = 1\) 于是蕴含 \(|\lambda|^2 \|b_{n+1}\|^2 = 1\)，而条件 \(\langle a_{n+1}|e_{n+1} \rangle > 0\) 蕴含 \(\lambda \langle a_{n+1}|b_{n+1} \rangle > 0\)；这就完全确定了 \(\lambda\)，于是我们证明了可以用归纳法确定唯一的标准正交族 \((e_n)_{n\in I}\)，使之满足命题的条件 1) 与 2)。

称序列 \((e_{n})_{n\in\mathbb{I}}\) 是由无关族 \((a_{n})_{n\in\mathbb{I}}\) 经标准正交化得到的。显然，由族 \((e_{n})\) 生成的向量子空间与由族 \((a_{n})\) 生成的向量子空间相同。特别地，若 \((a_{n})\) 是完全序列，则 \((e_{n})\) 也是，从而它是 E 的标准正交基；由此得到：

推论. --- 在每个满足第一可数公理的分离准希尔伯特空间 E 中，都存在可数的标准正交基。

若 E 满足第一可数公理，则 E 中存在完全序列，而且总能从此序列中抽取一个无关的完全族（A, II, § 7, No.~1, 定理 2）。

可以举出没有任何标准正交基的分离准希尔伯特空间的例子（V, p.~70, 练习 2）。

例. --- 设 I 为 R 的区间 \((-1,1)\)，E 为 I 上实值连续函数的向量空间。用 x 表示从 I 到 R 中的典范单射，把它看作 E 的一个元素。由 Stone-Weierstrass 定理，序列 \((x^{n})_{n\in\mathbf{N}}\) 在 E 中对一致收敛拓扑是完全的（GT, X, § 4, No.~2）。

把 E 看作实的分离准希尔伯特空间，其内积由下式给出

\[
\langle f | g \rangle = \int_ {- 1} ^ {1} f (t) g (t) d t.
\]

于是序列 \((x^{n})_{n\in \mathbb{N}}\) 在准希尔伯特空间 E 中是完全的。设 \((\Pi_n)_{n\in \mathbb{N}}\) 为由序列 \((x^n)_{n\in \mathbb{N}}\) 经标准正交化得到的序列。可以证明 \(\Pi_n = (n + \frac{1}{2})^{1 / 2}\mathbf{P}_n\)，其中 Legendre 多项式 \(\mathbf{P}_n\) 由下式定义

\[
\mathrm{P} _ {n} (x) = \frac {1}{2 ^ {n} n !} \left(\frac {d}{d x}\right) ^ {n} (x ^ {2} - 1) ^ {n}.
\]

命题 7. --- 在希尔伯特空间 E 中，两个标准正交基是等势的。

设 B 与 C 为 E 的两个标准正交基。当两集 B、C 之一为有限时，情形是平凡的，因为有限的标准正交基是空间的代数基。因此设 B 与 C 都是无限的。对每个 \(x \in B\)，设 \(C_{x}\) 为由 \(y \in C\) 中所有使 \(\langle x | y \rangle \neq 0\) 者组成的子集。集 \(C_{x}\) 是可数的（V, p.~21, 命题 4）。对每个 \(y \in C\)，存在 \(x \in B\) 使 \(y \in C_{x}\)，因为 B 是标准正交基且 \(y \neq 0\)；换言之，C 是当 x 取遍 B 时诸可数集 \(C_{x}\) 的并。于是 C 的基数小于 N × B 的基数，从而小于 B 的基数（S, III, §6, No.~3, 推论 4）；类似地，B 的基数小于 C 的基数；证明完毕。

希尔伯特空间 E 的任一标准正交基的基数称为 E 的希尔伯特维数。

推论 1. --- 给定希尔伯特空间 E 中的两个标准正交基，存在 E 的自同构把第一个基变为第二个基。

推论 2. --- 为使希尔伯特空间 \(\ell^2 (\mathbf{I})\) 与 \(\ell^2 (\mathbf{J})\) 同构，必要且充分的是 I 与 J 等势。

\section{希尔伯特空间的张量积}

\subsection{准希尔伯特空间的张量积}

设 \(E_{1}\) 与 \(E_{2}\) 为两个准希尔伯特空间，\(F = E_{1} \otimes E_{2}\) 为向量空间 \(E_{1}\) 与 \(E_{2}\) 的张量积。设 \(x_{1} \in E_{1}\) 且 \(x_{2} \in E_{2}\)；由于映射 \((y_{1}, y_{2}) \mapsto \langle x_{1}|y_{1} \rangle \langle x_{2}|y_{2} \rangle\)（从 \(E_{1} \times E_{2}\) 到 K 中）是双线性的，故存在线性型 \(\phi_{x_{1},x_{2}}\)（在 \(E_{1} \otimes E_{2}\) 上），使得

\[
\phi_ {x _ {1}, x _ {2}} (y _ {1} \otimes y _ {2}) = \langle x _ {1} | y _ {1} \rangle \langle x _ {2} | y _ {2} \rangle\tag{1}
\]

其中 \(y_{1} \in \mathrm{E}_{1}\)，\(y_{2} \in \mathrm{E}_{2}\)。设 \(z \in \mathrm{F}\)。映射 \((x_{1}, x_{2}) \mapsto \overline{\phi_{x_{1},x_{2}}(z)}\)（从 \(\mathrm{E}_1 \times \mathrm{E}_2\) 到 \(\mathbf{K}\) 中）是双线性的；把 \(z\) 写成形式 \(z = \sum_{i=1}^{n} y_{i,1} \otimes y_{i,2}\)（其中 \(y_{i,1} \in \mathrm{E}_1\)、\(y_{i,2} \in \mathrm{E}_2\)，\(1 \leqslant i \leqslant n\)）即可看出这一点。于是存在线性型 \(\psi_z\)（在 \(\mathrm{F} = \mathrm{E}_1 \otimes \mathrm{E}_2\) 上），使得

\[
\psi_ {z} (x _ {1} \otimes x _ {2}) = \overline {{\phi_ {x _ {1} , x _ {2}} (z)}} \quad (x _ {1} \in \mathrm{E} _ {1}, x _ {2} \in \mathrm{E} _ {2}).\tag{2}
\]

我们令 \(\Phi(z, t) = \psi_z(t)\)（对 F 中的 \(z, t\)）。立即可以看出 \(\Phi\) 是 \(\mathbf{E}_1 \otimes \mathbf{E}_2\) 上的一个半双线性型，它由下式刻画

\[
\Phi (x _ {1} \otimes x _ {2}, y _ {1} \otimes y _ {2}) = \langle x _ {1} | y _ {1} \rangle \langle x _ {2} | y _ {2} \rangle\tag{3}
\]

(cf.~A, IX, § 1, No.~11).

命题 1. --- 半双线性型 \(\Phi\)（在 \(\mathrm{E}_1 \otimes \mathrm{E}_2\) 上）是埃尔米特型且是正的，因而它给 \(\mathrm{E}_1 \otimes \mathrm{E}_2\) 赋予了准希尔伯特空间的结构。若 \(\mathrm{E}_1\) 与 \(\mathrm{E}_2\) 是分离的，则此空间是分离的。

公式 \(\Phi(z, t) = \overline{\Phi(t, z)}\) 当 \(z = x_1 \otimes x_2\) 且 \(t = y_1 \otimes y_2\) 时由 (3) 得出。一般情形由线性性得到，因而 \(\Phi\) 是埃尔米特型。

设 \(\mathbf{E}_1\) 与 \(\mathbf{E}_2\) 是分离的；我们将证明埃尔米特型 \(\Phi\) 是正的且非退化。设 \(z = \sum_{i=1}^{n} x_i \otimes y_i\) 为 \(\mathrm{F} = \mathrm{E}_1 \otimes \mathrm{E}_2\) 的非零元素。设 \((e_1, ..., e_n)\) 为 \(\mathbf{E}_1\) 中由 \(x_1, ..., x_n\) 生成的子空间的标准正交基（V, p.~23, 推论 1）。存在元素 \(f_1, ..., f_m\)（在 \(\mathbf{E}_2\) 中，不全为零），使得 \(z = \sum_{i=1}^{n} e_i \otimes f_i\)，因而

\[
\begin{array}{l} \Phi (z, z) = \sum_ {i, j = 1} ^ {m} \Phi (e _ {i} \otimes f _ {i}, e _ {j} \otimes f _ {j}) \\ = \sum_ {i, j} \langle e _ {i} | e _ {j} \rangle \langle f _ {i} | f _ {j} \rangle = \sum_ {i = 1} ^ {m} \| f _ {i} \| ^ {2} > 0. \end{array}
\]

对一般情形，我们现在来证明 \(\Phi\) 是正的。设 \(\tilde{E}_{i}\) 为与 \(E_{i}\) 相伴的分离准希尔伯特空间，\(\pi_{i}\) 为从 \(E_{i}\) 到 \(\tilde{E}_{i} (i = 1, 2)\) 上的典范映射。令 \(\pi = \pi_{1} \otimes \pi_{2}\)。设 \(\tilde{\Phi}\) 为 \(\tilde{E}_{1} \otimes \tilde{E}_{2}\) 上按与 \(\Phi\) 相同的方式构造的埃尔米特型。我们有

\[
\Phi (z, t) = \Phi (\pi (z), \pi (t)) \quad (z \in \mathrm{F}, t \in \mathrm{F}),
\]

而由于 \(\tilde{\Phi}\) 是正的，\(\Phi\) 也是正的。

证毕。

命题 1 中定义的准希尔伯特空间称为准希尔伯特空间 \(E_{1}\) 与 \(E_{2}\) 的张量积，记作 \(E_{1} \otimes E_{2}\)。此后我们将用 \(\langle z|t\rangle\) 表示 \(\Phi(z,t)\)，于是依定义

\[
\langle x _ {1} \otimes x _ {2} | y _ {1} \otimes y _ {2} \rangle = \langle x _ {1} | y _ {1} \rangle \langle x _ {2} | y _ {2} \rangle ;\tag{4}
\]

我们也用 \(\| z\| _2\) 或 \(\| z\|\) 来记 \(\langle z|z\rangle^{1 / 2}\)。由 (4) 得

\[
\left\| x _ {1} \otimes x _ {2} \right\| _ {2} = \left\| x _ {1} \right\|. \left\| x _ {2} \right\|,\tag{5}
\]

于是双线性映射 \((x_{1}, x_{2}) \mapsto x_{1} \otimes x_{2}\)（从 \(\mathrm{E}_1 \times \mathrm{E}_2\) 到 \(\mathrm{E}_1 \otimes \mathrm{E}_2\) 中）是连续的。

对 i = 1, 2，设 \(F_{i}\) 为 \(E_{i}\) 的向量子空间，赋以诱导的准希尔伯特结构。于是 \(F_{1} \otimes F_{2}\) 可与 \(E_{1} \otimes E_{2}\) 的一个向量子空间等同（A, II, § 7, No.~7）。公式 (4) 表明，\(F_{1} \otimes F_{2}\) 赋以由 \(E_{1} \otimes_{2} E_{2}\) 的结构诱导的准希尔伯特空间结构后，恰好就是 \(F_{1} \otimes_{2} F_{2}\)。此后我们将把 \(F_{1} \otimes_{2} F_{2}\) 与 \(E_{1} \otimes_{2} E_{2}\) 的准希尔伯特子空间等同起来。

命题 2. --- 对 \(i = 1, 2\)，设 \(\mathrm{E}_{i}\) 与 \(\mathrm{F}_{i}\) 为两个分离的准希尔伯特空间，并设 \(u_{i} \in \mathcal{L}(\mathrm{E}_{i}; \mathrm{F}_{i})\)。线性映射 \(u_{1} \otimes u_{2}\)（从 \(\mathrm{E}_{1} \otimes_{2} \mathrm{E}_{2}\) 到 \(\mathrm{F}_{1} \otimes_{2} \mathrm{F}_{2}\) 中）是连续的，并且我们有

\[
\left\| u _ {1} \otimes u _ {2} \right\| = \left\| u _ {1} \right\|. \left\| u _ {2} \right\|.
\]

考虑 \(\mathbf{E}_1\) 上由下式给出的正埃尔米特型

\[
f \left(x _ {1}, y _ {1}\right) = \| u _ {1} \| ^ {2} \left\langle x _ {1} \mid y _ {1} \right\rangle - \left\langle u _ {1} \left(x _ {1}\right) \mid u _ {1} \left(y _ {1}\right) \right\rangle .
\]

由命题 1（V, p.~25），存在正埃尔米特型 \(\Phi\)（在 \(E_1 \otimes E_2\) 上），使得

\[
\begin{array}{r l} \Phi (x _ {1} \otimes x _ {2}, y _ {1} \otimes y _ {2}) & = f (x _ {1}, y _ {1}) \langle x _ {2} | y _ {2} \rangle = \\ & = \| u _ {1} \| ^ {2} \langle x _ {1} \otimes x _ {2} | y _ {1} \otimes y _ {2} \rangle - \langle (u _ {1} \otimes 1) (x _ {1} \otimes x _ {2}) | (u _ {1} \otimes 1) (y _ {1} \otimes y _ {2}) \rangle \end{array}
\]

其中 \(x_{1}, y_{1}\) 属于 \(E_{1}\)，\(x_{2}, y_{2}\) 属于 \(E_{2}\)。由线性性得

\[
\Phi (z, t) = \| u _ {1} \| ^ {2} \langle z | t \rangle - \langle (u _ {1} \otimes 1) (z) | (u _ {1} \otimes 1) (t) \rangle
\]

其中 \(z, t\) 属于 \(\mathrm{E}_1 \otimes \mathrm{E}_2\)。由于 \(\Phi\) 是正的，得 \(\Phi(z, z) \geqslant 0\)，即 \(\| (u_1 \otimes 1).z \|_2\) \(\leqslant \| u_1\| \cdot \| z\| _2\)（对 \(z\in \mathrm{E}_1\otimes_2\mathrm{E}_2\)），亦即 \(\| u_1\otimes 1\| \leqslant \| u_1\|\)。同理可证不等式 \(\| 1\otimes u_2\| \leqslant \| u_2\|\)；又由于 \(u_{1}\otimes u_{2} = (u_{1}\otimes 1)\circ (1\otimes u_{2})\)，我们得到

\[
\left\| u _ {1} \otimes u _ {2} \right\| \leqslant \left\| u _ {1} \right\|. \left\| u _ {2} \right\|.
\]

另一方面，

\[
\begin{array}{l} \| u _ {1} \| \cdot \| u _ {2} \| = \sup _ {\| x _ {1} \| \leqslant 1, \| x _ {2} \| \leqslant 1} \| u _ {1} (x _ {1}) \| \cdot \| u _ {2} (x _ {2}) \| \\ = \sup _ {\| x _ {1} \| \leqslant 1, \| x _ {2} \| \leqslant 1} \| (u _ {1} \otimes u _ {2}) (x _ {1} \otimes x _ {2}) \| _ {2} \leqslant \| u _ {1} \otimes u _ {2} \|. \end{array}
\]

这样就完成了命题 2 的证明。

证毕。

设 \(E_1, \ldots, E_n\) 为准希尔伯特空间（\(n \geqslant 2\)）。我们用归纳法以下式定义张量积 \(E_1 \otimes_2 \ldots \otimes_2 E_n\)（也记作 \(\bigotimes_{i=1}^{n} E_i\)）

\[
\mathrm{E} _ {1} \otimes_ {2} \dots \otimes_ {2} \mathrm{E} _ {n} = (\mathrm{E} _ {1} \otimes_ {2} \dots \otimes_ {2} \mathrm{E} _ {n - 1}) \otimes_ {2} \mathrm{E} _ {n}.
\]

于是，由内积的定义，我们有

\[
\left\langle x _ {1} \otimes \dots \otimes x _ {n} \mid y _ {1} \otimes \dots \otimes y _ {n} \right\rangle = \prod_ {i = 1} ^ {n} \left\langle x _ {i} \mid y _ {i} \right\rangle ,\tag{6}
\]

特别地 \(^{1}\)

\[
\left\| x _ {1} \otimes \dots \otimes x _ {n} \right\| _ {2} = \left\| x _ {1} \right\| \dots \left\| x _ {n} \right\|,\tag{7}
\]

其中 \(x_{i},y_{i}\) 属于 \(\mathrm{E}_i\)（\(1\leqslant i\leqslant n\)）。若诸 \(\mathrm{E}_i\) 都是分离的，则 \(\mathrm{E}_1\otimes_2\dots \otimes_2\mathrm{E}_n\) 也是分离的。

设 \(F_{1}, \ldots, F_{n}\) 为准希尔伯特空间，且 \(u_{i} \in \mathcal{L}(E_{i}; F_{i})\)（\(1 \leqslant i \leqslant n\)）。对 n 用归纳法，由命题 2 推出 \(u_{1} \otimes \ldots \otimes u_{n}\) 是从 \(E_{1} \otimes_{2} \ldots \otimes_{2} E_{n}\) 到 \(F_{1} \otimes_{2} \ldots \otimes_{2} F_{n}\) 中的连续线性映射，并且

\[
\left\| u _ {1} \otimes \dots \otimes u _ {n} \right\| = \left\| u _ {1} \right\| \dots \left\| u _ {n} \right\|.\tag{8}
\]

设 \(\sigma \in \mathfrak{S}_n\) 为集合 \(\{1, 2, \dots, n\}\) 的一个置换。由 (6)，线性映射 \(p_{\sigma}\)（从 \(\mathrm{E}_1 \otimes_2 \dots \otimes_2 \mathrm{E}_n\) 到 \(\mathrm{E}_{\sigma^{-1}(1)} \otimes_2 \dots \otimes_2 \mathrm{E}_{\sigma^{-1}(n)}\) 上）由下式刻画

\[
p _ {\sigma} \left(x _ {1} \otimes \dots \otimes x _ {n}\right) = x _ {\sigma^ {- 1} (1)} \otimes \dots \otimes x _ {\sigma^ {- 1} (n)}\tag{9}
\]

是准希尔伯特空间同构（« 张量积的交换性 »）。

类似地，考虑把 \(\{1, 2, \dots , n\}\) 分成 \(m\) 个相邻区间 \(\mathrm{I}_{1}, \dots , \mathrm{I}_{m}\) 的一个划分，其中 \(\mathrm{I}_{k} = [a_{k}, a_{k + 1} - 1]\)（\(1 \leqslant k \leqslant m\)）。令

\[
\mathrm{F} _ {k} = \bigotimes_ {i = a _ {k}} ^ {a _ {k + 1} - 1} \mathrm{E} _ {i} (1 \leqslant k \leqslant m).
\]

\(^{1}\) 这里我们同样令 \(\|z\|_{2} = \langle z|z\rangle^{1/2}\)（z 在 \(E_{1} \otimes_{2} \ldots \otimes_{2} E_{n}\) 中）。

从 \(F_{1} \otimes \ldots \otimes F_{m}\) 到 \(E_{1} \otimes \ldots \otimes E_{n}\) 上的、把 \(\bigotimes_{k=1}^{m} \bigotimes_{i=a_{k}}^{a_{k+1}-1} x_{i}\) 变为 \(x_{1} \otimes \ldots \otimes x_{n}\) 的典范同构（A, II, §3, No.~9）是准希尔伯特空间同构（« 张量积的结合性 »）。

\subsection{希尔伯特空间的希尔伯特张量积}

定义 1. --- 设 \(E_{1}, \ldots, E_{n}\) 为希尔伯特空间。分离准希尔伯特空间 \(E_{1} \otimes_{2} \ldots \otimes_{2} E_{n}\) 的完备化称为诸 \(E_{i}\) 的希尔伯特张量积，记作 \(E_{1} \hat{\otimes}_{2} \ldots \hat{\otimes}_{2} E_{n}\)（或 \(\hat{\otimes}_{2} E_{i}\)）。

设 \(F_{1}, \ldots, F_{n}\) 为希尔伯特空间，且 \(u_{i} \in \mathcal{L}(E_{i}, F_{i})\)（\(1 \leqslant i \leqslant n\)）。于是连续线性映射 \(u_{1} \otimes \ldots \otimes u_{n}\) 可延拓为连续线性映射 \(u_{1} \hat{\otimes}_{2} \ldots \hat{\otimes}_{2} u_{n}\)（从 \(E_{1} \hat{\otimes}_{2} \ldots \hat{\otimes}_{2} E_{n}\) 到 \(F_{1} \hat{\otimes}_{2} \ldots \hat{\otimes}_{2} F_{n}\) 中）。我们有

\[
\| u _ {1} \hat {\otimes} _ {2} \dots \hat {\otimes} _ {2} u _ {n} \| = \| u _ {1} \| \dots \| u _ {n} \|\tag{10}
\]

这由 V, p.~27 的公式 (8) 得到。此外，若以 \(1_{E}\) 表示任意希尔伯特空间 E 的恒等映射，则有

\[
1 _ {\mathrm{E} _ {1}} \hat {\otimes} _ {2} \dots \hat {\otimes} _ {2} 1 _ {\mathrm{E} _ {n}} = 1 _ {\mathrm{E}} \quad \text { with } \quad \mathrm{E} = \mathrm{E} _ {1} \hat {\otimes} _ {2} \dots \hat {\otimes} _ {2} \mathrm{E} _ {n}.\tag{11}
\]

最后，若 \(\mathbf{G}_1, \ldots, \mathbf{G}_n\) 为希尔伯特空间且 \(v_i \in \mathcal{L}(\mathbf{F}_i; \mathbf{G}_i)\)（\(1 \leqslant i \leqslant n\)），则得到

\[
\left(v _ {1} \circ u _ {1}\right) \hat {\otimes} _ {2} \dots \hat {\otimes} _ {2} \left(v _ {n} \circ u _ {n}\right) = \left(v _ {1} \hat {\otimes} _ {2} \dots \hat {\otimes} _ {2} v _ {n}\right) \circ \left(u _ {1} \hat {\otimes} _ {2} \dots \hat {\otimes} _ {2} u _ {n}\right).\tag{12}
\]

与上面对准希尔伯特空间所述的内容相类比，希尔伯特张量积的«交换性»与«结合性»如何表述，留给读者自己去完成。

注. --- 设 \(E_{1}, \ldots, E_{n}\) 为分离的准希尔伯特空间，\(\hat{E}_{1}, \ldots, \hat{E}_{n}\) 为它们各自的完备化。于是 \(E_{1} \otimes_{2} \ldots \otimes_{2} E_{n}\) 是 \(\hat{E}_{1} \otimes_{2} \ldots \otimes_{2} \hat{E}_{n}\) 的准希尔伯特子空间。由于映射 \((x_{1}, \ldots, x_{n}) \mapsto x_{1} \otimes \ldots \otimes x_{n}\)（从 \(\hat{E}_{1} \times \cdots \times \hat{E}_{n}\) 到 \(\hat{E}_{1} \otimes_{2} \ldots \otimes_{2} \hat{E}_{n}\) 中）是连续的，故 \(E_{1} \otimes_{2} \ldots \otimes_{2} E_{n}\) 在 \(\hat{E}_{1} \otimes_{2} \ldots \otimes_{2} \hat{E}_{n}\) 中稠密。更有甚者，\(E_{1} \otimes_{2} \ldots \otimes_{2} E_{n}\) 的完备化恰好就是希尔伯特空间 \(\hat{E}_{1} \otimes_{2} \ldots \otimes_{2} \hat{E}_{n}\)。这一完备化有时就简记作 \(E_{1} \otimes_{2} \ldots \otimes_{2} E_{n}\)（或 \(\bigotimes_{1 \leq i \leq n} E_{i}\)）。

命题 3. --- 设 \(E_{1}, \ldots, E_{n}\) 为希尔伯特空间。假定对 \(1 \leqslant i \leqslant n\)，空间 \(E_{i}\) 是闭向量子空间族 \((\mathrm{E}_{i,\alpha})_{\alpha \in \mathbf{A}(i)}\) 的希尔伯特和。于是 \(E_{1} \otimes_{2} \ldots \otimes_{2} E_{n}\) 是子空间族 \(E_{1,\alpha_{1}} \hat{\otimes}_{2} \ldots \hat{\otimes}_{2} E_{n,\alpha_{n}}\) 的希尔伯特和，其中 \((\alpha_{1}, \ldots, \alpha_{n})\) 取遍 \(A(1) \times \cdots \times A(n)\)。

由 V, p.~27 的公式 (6)，子空间 \(\mathrm{E}_{1,\alpha_1}\hat{\otimes}_2\ldots \hat{\otimes}_2\mathrm{E}_{n,\alpha_n}\)（\(\mathrm{E}_1\hat{\otimes}_2\ldots \hat{\otimes}_2\mathrm{E}_n\) 中的）两两正交。对每个整数 \(i\)（介于 1 与 \(n\) 之间），集合 \(\bigcup_{\alpha \in \mathbf{A}(i)}\mathrm{E}_{i,\alpha}\) 在 \(\mathrm{E}_i\) 中是完全的，且多线性映射 \((x_{1},\dots,x_{n})\mapsto x_{1}\otimes \dots \otimes x_{n}\) 是连续的。由此推出诸子空间 \(\mathrm{E}_{1,\alpha_1}\hat{\otimes}_2\ldots \hat{\otimes}_2\mathrm{E}_{n,\alpha_n}\) 的并是完全的，从而命题 3 成立。

推论 1. --- 对 \(1 \leqslant i \leqslant n\)，设 \((e_{i,\alpha})_{\alpha \in \mathbf{A}(i)}\) 为 \(\mathbf{E}_i\) 的标准正交基。于是向量族 \(e_{1,\alpha_1} \otimes \ldots \otimes e_{n,\alpha_n}\)（当 \((\alpha_1, \ldots, \alpha_n)\) 取遍 \(\mathbf{A}(1) \times \cdots \times \mathbf{A}(n)\) 时）是 \(\mathbf{E}_1 \hat{\otimes}_2 \ldots \hat{\otimes}_2 \mathbf{E}_n\) 的标准正交基。

推论 2. --- 设 \(E_{1}\) 与 \(E_{2}\) 为两个希尔伯特空间，\((e_{i})_{i\in I}\) 为 \(E_{1}\) 的标准正交基。设 \((y_{i})_{i\in I}\) 为 \(E_{2}\) 中元素的一个族，满足 \(\sum_{i\in I}\|y_{i}\|^{2}<+\infty\)。于是族 \((e_{i}\otimes y_{i})_{i\in I}\) 在 \(E_{1}\hat{\otimes}_{2}E_{2}\) 中可和；此外，\(E_{1}\hat{\otimes}_{2}E_{2}\) 的每个元素都可以唯一地写成 \(\sum_{i\in I}e_{i}\otimes y_{i}\) 的形式，其中 \(\sum_{i\in I}\|y_{i}\|^{2}<+\infty\)。

设 \(F_{i}\) 为 \(E_{1}\) 中由 \(e_{i} (i \in I)\) 生成的直线。于是 \(E_{1}\) 是子空间族 \((F_{i})_{i \in I}\) 的希尔伯特和。由命题 3，空间 \(E_{1} \hat{\otimes} E_{2}\) 是子空间族 \((F_{i} \hat{\otimes}_{2} E_{2})_{i \in I}\) 的希尔伯特和，由此即得推论 2。

例. --- 1) 由推论 1，空间 \(\ell^2 (\mathrm{I})\hat{\otimes}_2\ell^2 (\mathrm{J})\) 典范同构于 \(\ell^2 (\mathrm{I}\times \mathrm{J})\)，张量积 \(\mathbf{x}\otimes \mathbf{y}\)（其中 \(\mathbf{x} = (x_i)_{i\in \mathrm{I}}\)，\(\mathbf{y} = (y_j)_{j\in \mathrm{J}}\)）可与族 \((x_iy_j)_{i\in \mathrm{I},j\in \mathrm{J}}\) 等同。类似地，由推论 2，\(\ell^2 (\mathrm{I})\hat{\otimes}_2\mathrm{E}\) 可与 \(\ell_{\mathrm{E}}^{2}(\mathrm{I})\) 等同，使得 \((x_{i})_{i\in \mathrm{I}}\otimes y = (x_{i}y)_{i\in \mathrm{I}}\)（对希尔伯特空间 E 中的每个 \(\mathbf{y}\)）。

* 2) 设 X 为分离的拓扑空间，\(\mu\) 为 X 上的一个正测度。设 E 为希尔伯特空间。我们可以用典范的方式把 \(L^2(X, \mu) \otimes_2 E\) 与 \(L_E^2(X, \mu)\) 等同起来：若 \(f\) 是 X 上平方可积数值函数 \(f\) 的类，且 \(a\) 属于 E，则 \(f \otimes a\) 就是取值于 E 中的函数 \(x \mapsto f(x).a\) 的类。

设 Y 为分离的拓扑空间，v 为 Y 上的一个正测度。以类似的方式，我们可以把希尔伯特空间 \(\mathrm{L}^{2}(\mathrm{X},\mu)\hat{\otimes}_{2}\mathrm{L}^{2}(\mathrm{Y},v)\) 与 \(\mathrm{L}^{2}(\mathrm{X}\times\mathrm{Y},\mu\otimes v)\) 等同起来；这时 \(\dot{f}\otimes\dot{g}\) 可与函数 \((x,y)\mapsto f(x)g(y)\)（在 \(X\times Y\) 上）的类等同。

\subsection{对称希尔伯特幂}

设 E 为希尔伯特空间，n 为正整数。以 \(\hat{\mathbf{T}}^{n}(\mathrm{E})\) 或 \(E^{\otimes_{n}}\) 表示 n 个都等于 E 的希尔伯特空间的张量积。换言之，\(\hat{\mathbf{T}}^{n}(\mathrm{E})\) 是空间 \(\mathbf{T}^{n}(\mathrm{E}) = \mathrm{E} \otimes \ldots \otimes \mathrm{E}\)（n 个因子）关于由下式定义的分离准希尔伯特空间结构的完备化

\[
\left\langle x _ {1} \otimes \dots \otimes x _ {n} \mid y _ {1} \otimes \dots \otimes y _ {n} \right\rangle = \prod_ {i = 1} ^ {n} \left\langle x _ {i} \mid y _ {i} \right\rangle .
\]

若 \((e_i)_{i\in I}\) 是 \(E\) 的标准正交基，则向量族 \(e_{i_1} \otimes \ldots \otimes e_{i_n}\)（\(i_1, \ldots, i_n\) 取遍 \(I\)）是 \(\hat{\mathbf{T}}^n(E)\) 的标准正交基（V, p.~29, 推论 1）。我们得到 \(\hat{\mathbf{T}}^\circ(E) = K\)。

设 \(\sigma \in \mathfrak{S}_n\) 为集合 \(\{1, 2, ..., n\}\) 的一个置换。由 V, p.~27，存在自同构 \(p_{\sigma}\)（\(\hat{\mathbf{T}}^n(\mathbf{E})\) 的），它由下式刻画

\[
p _ {\sigma} (x _ {1} \otimes \dots \otimes x _ {n}) = x _ {\sigma^ {- 1} (1)} \otimes \dots \otimes x _ {\sigma^ {- 1} (n)}.
\]

我们有 \(p_{\sigma \tau} = p_{\sigma}p_{\tau}\)（\(\sigma, \tau\) 属于 \(\mathfrak{S}_n\)），因而自同态 \(\Pi_n = \frac{1}{n!}\sum_{\sigma \in \mathfrak{S}_n}p_\sigma\)（向量空间 \(\hat{\mathbf{T}}^n (\mathrm{E})\) 的）是到在 \(\mathfrak{S}_n\) 下保持不变的所有元素组成的子空间上的正交投影算子。然而（A, III, § 6, No.~3），\(\Pi_n\) 把« 代数 »张量积 \(\mathbf{T}^n (\mathrm{E})\) 映到子空间 \(\mathbf{T}\mathbf{S}^n (\mathrm{E})\)（由所有 \(n\) 阶对称张量组成的）上。换言之，\(\Pi_n\) 的像是空间 \(\mathbf{T}\mathbf{S}^n (\mathrm{E})\)（赋以由 \(\mathbf{T}^n (\mathrm{E})\) 的内积诱导的内积）的完备化；这一完备化将记作 \(\widehat{\mathbf{T}\mathbf{S}}^n (\mathrm{E})\)。

设 \(\mathbf{S}^n (\mathrm{E})\) 为向量空间 E 的 \(n\) 次对称幂（A, III, § 6, No.~1）。从 \(\mathbf{T}^n (\mathrm{E})\) 到 \(\mathbf{S}^n (\mathrm{E})\) 上的典范映射通过限制定义了一个同构 \(\lambda_{n}\)（从 \(\mathbf{TS}^n (\mathrm{E})\) 到 \(\mathbf{S}^n (\mathrm{E})\) 上）。我们立即可以验证，其逆同构由下式给出

\[
\mu_ {n} (x _ {1} \dots x _ {n}) = \Pi_ {n} (x _ {1} \otimes \dots \otimes x _ {n}) = \frac {1}{n !} \sum_ {\sigma \in \mathfrak {S} _ {n}} x _ {\sigma^ {- 1} (1)} \otimes \dots \otimes x _ {\sigma^ {- 1} (n)}\tag{13}
\]

其中 \(x_{1},\ldots ,x_{n}\) 属于 E。

我们通过令下式成立而在 \(\mathbf{S}^{n}(\mathrm{E})\) 上定义一个分离准希尔伯特空间结构

\[
\langle u | v \rangle = n! \left\langle \mu_ {n} (u) | \mu_ {n} (v) \right\rangle .\tag{14}
\]

于是我们有（试与 A, III, § 11, No.~5 的公式 (29) 比较）

\[
\left\langle x _ {1} \dots x _ {n} \mid y _ {1} \dots y _ {n} \right\rangle = \sum_ {\sigma \in \mathfrak {S} _ {n}} \prod_ {i = 1} ^ {n} \left\langle x _ {i} \mid y _ {\sigma (i)} \right\rangle ,\tag{15}
\]

特别地

\[
\langle x ^ {n} | y ^ {n} \rangle = n! \langle x | y \rangle^ {n}.\tag{16}
\]

设 \(\hat{\mathbf{S}}^n (\mathrm{E})\) 表示准希尔伯特空间 \(\mathbf{S}^n (\mathrm{E})\) 的完备化，\(\hat{\mathbf{S}} (\mathrm{E})\) 表示诸希尔伯特空间 \(\hat{\mathbf{S}}^n (\mathrm{E})\) 的外希尔伯特和。我们可以证明（V, p.~73, 练习 1），除非 E 恰为 0，代数 \(\mathbf{S}(\mathrm{E})\) 中的乘法不可能连续地延拓到 \(\mathbf{S}(\mathrm{E})\) 上。

命题 4. --- 设 \((e_{i})_{i\in I}\) 为希尔伯特空间 E 的标准正交基。对每个 \(\alpha\)（属于 \(\mathbf{N}^{(1)}\)），令

\[
z _ {\alpha} = \prod_ {i \in I} e _ {i} ^ {\alpha_ {i}} / (\alpha_ {i}!) ^ {1 / 2}.\tag{17}
\]

于是 \((z_{\alpha})_{\alpha \in \mathbf{N}^{(1)}}\) 是 \(\hat{\mathbf{S}} (\mathrm{E})\) 的标准正交基。

设 \(E_{0}\) 为 E 中由向量 \(e_{i}\)（i 取遍 I）生成的向量子空间。于是诸 \(z_{\alpha}\) 构成向量空间 \(\mathbf{S}(E_{0})\) 的基（A, III, § 6, No.~6）。但 \(E_{0}\) 在 E 中稠密，且多线性映射 \((x_{1},...,x_{n})\mapsto x_{1}\ldots x_{n}\)（从 \(E\times\cdots\times E\) 到 \(\mathbf{S}(E)\) 中）对一切 \(n\geqslant1\) 是连续的；因此 \(\mathbf{S}(E_{0})\) 在 \(\mathbf{S}(E)\) 中稠密。现在只需证明诸 \(z_{\alpha}\) 组成的族是标准正交的。首先注意，\(\hat{\mathbf{S}}^{n}(\mathrm{E})\) 与 \(\hat{\mathbf{S}}^{m}(\mathrm{E})\)（当 \(n \neq m\) 时）正交。于是只需证明公式

\[
\langle z _ {\alpha} | z _ {\beta} \rangle = \left\{ \begin{array}{l l} 1 & \text {if} \quad \alpha = \beta \\ 0 & \text {if} \quad \alpha \neq \beta \end{array} \right.
\]

其中 \(|\alpha| = \sum_{i \in I} \alpha_i\) 与 \(|\beta| = \sum_{i \in I} \beta_i\) 等于同一个整数 \(n\)。

考虑一个划分 \((\mathbf{P}_i)_{i\in I}\)（集合 \(\{1,2,\dots,n\}\) 的划分），使得 \(\operatorname{Card} \mathbf{P}_i = \alpha_i\) 对所有 \(i\in I\) 成立。令 \(x_k = e_i\)（若 \(k\) 属于 \(\mathbf{P}_i\)），于是 \(x_1\ldots x_n = \prod_{i\in I}e_i^{\alpha_i}\)。类似地，我们定义 \((\mathbf{Q}_i)_{i\in I}\) 与 \(y_k\)，使得 \(\operatorname{Card} \mathbf{Q}_i = \beta_i\) 且 \(y_1\ldots y_n = \prod_{i\in I}e_i^{\beta_i}\)。由于诸 \(e_i\) 两两正交，我们有 \(\langle x_k|y_{\sigma(k)}\rangle = 0\)，除非存在指标 \(i\in I\) 使 \(k\in \mathbf{P}_i\) 且 \(\sigma (k)\in \mathbf{Q}_i\)。于是由公式 (15)，我们有 \(\langle x_1\ldots x_n|y_1\ldots y_n\rangle = 0\)，除非存在置换 \(\sigma \in \mathfrak{S}_n\) 使 \(\sigma (\mathbf{P}_i) = \mathbf{Q}_i\) 对所有 \(i\in I\) 成立（这蕴含 \(\alpha = \beta\)）。于是 \(\langle z_{\alpha}|z_{\beta}\rangle = 0\)（当 \(\alpha \neq \beta\) 时）。同样的论证证明 \(\| x_1\ldots x_n\|^2\) 等于 \(\sigma \in \mathfrak{S}_n\) 中使 \(\sigma (\mathbf{P}_i) = \mathbf{P}_i\) 对所有 \(i\in I\) 成立者的个数，从而等于 \(\prod_{i=1}^{\infty}\alpha_i!\)。由此得 \(\| z_\alpha \| = 1\)，命题得证。

推论. --- 假定希尔伯特空间 E 是正交子空间 M 与 N 的直和。从 \(\mathbf{S}(\mathbf{M}) \otimes \mathbf{S}(\mathbf{N})\) 到 \(\mathbf{S}(\dot{\mathbf{E}})\) 上的典范同构 g（A, III, §6, No.~6）可唯一地延拓为从 \(\hat{\mathbf{S}}(\mathbf{M}) \otimes_{2} \hat{\mathbf{S}}(\mathbf{N})\) 到 \(\hat{\mathbf{S}}(\mathbf{E})\) 上的希尔伯特空间同构 h。

设 \((e_{i})_{i\in I}\)（相应地 \((f_{j})_{j\in J}\)）为希尔伯特空间 M（相应地 N）的标准正交基，并设 \(M_{0}\)（相应地 \(N_{0}\)）为 E 中由向量 \(e_{i}\)（相应地 \(f_{j}\)）生成的向量子空间。令 \(E_{0}=M_{0}+N_{0}\)，并以 \(g_{0}\) 表示从 \(S(M_{0})\otimes S(N_{0})\) 到 \(S(E_{0})\) 上的典范同构。令

\[
z _ {\alpha} = \prod_ {i \in \mathbf {I}} e _ {i} ^ {\alpha_ {i}} / (\alpha_ {i}!) ^ {1 / 2}, t _ {\beta} = \prod_ {j \in \mathbf {J}} f _ {j} ^ {\beta _ {j}} / (\beta_ {j}!) ^ {1 / 2}
\]

其中 \(\alpha \in \mathbf{N}^{(I)}\)，\(\beta \in \mathbf{N}^{(J)}\)。由命题 4，我们这样就定义了标准正交基 \((z_{\alpha})_{\alpha \in \mathbf{N}^{(I)}}\)（\(\hat{\mathbf{S}}(\mathbf{M})\) 的）、\((t_{\beta})_{\beta \in \mathbf{N}^{(J)}}\)（\(\hat{\mathbf{S}}(\mathbf{N})\) 的）以及 \((z_{\alpha}t_{\beta})_{\alpha \in \mathbf{N}^{(I)},\beta \in \mathbf{N}^{(J)}}\)（\(\hat{\mathbf{S}}(\mathbf{E})\) 的）。由于 \(z_{\alpha}t_{\beta} = g_0(z_{\alpha}\otimes t_{\beta})\)，且诸元素 \(z_{\alpha}\otimes t_{\beta}\) 构成 \(\hat{\mathbf{S}}(\mathbf{M})\hat{\otimes}_2\hat{\mathbf{S}}(\mathbf{N})\) 的标准正交基（V, p.~29, 推论 1），可见 \(g_0\) 可延拓为希尔伯特空间同构 \(h:\hat{\mathbf{S}}(\mathbf{M})\hat{\otimes}_2\hat{\mathbf{S}}(\mathbf{N})\to \hat{\mathbf{S}}(\mathbf{E})\)。根据构造，我们有

\[
h (x _ {1} \dots x _ {m} \otimes y _ {1} \dots y _ {n}) = x _ {1} \dots x _ {m} y _ {1} \dots y _ {n}
\]

其中向量 \(x_{1}, \ldots, x_{m}\) 属于 \(M_{0}\)，向量 \(y_{1}, \ldots, y_{n}\) 属于 \(N_{0}\)。由连续性，同一关系式对 M 中的向量 \(x_{1}, \ldots, x_{m}\) 与 N 中的向量 \(y_{1}, \ldots, y_{n}\) 也成立；换言之，h 延拓了 g。h 的唯一性是显然的。

设 E 与 F 为两个希尔伯特空间，\(u \in \mathcal{L}(E; F)\)。线性映射 \(\hat{\mathbf{T}}^{n}(u) = u \otimes_{2} \ldots \otimes_{2} u\)（n 个因子）是从 \(\hat{\mathbf{T}}^{n}(E)\) 到 \(\hat{\mathbf{T}}^{n}(F)\) 中的连续映射，其范数为 \(\|u\|^{n}\)（V, p.~28, 公式 (10)）。此外，V, p.~30 的公式 (13) 与 (14) 表明，存在一个同构 \(\phi_{n,\mathrm{E}}\)（从 \(\hat{\mathbf{S}}^n (\mathrm{E})\) 到子空间 \(\widehat{\mathbf{T}\mathbf{S}}^n (\mathrm{E})\)（\(\hat{\mathbf{T}}^n (\mathrm{E})\) 的）上）且仅存在一个，使得

\[
\phi_ {n, \mathrm{E}} (x _ {1} \dots x _ {n}) = \frac {1}{(n !) ^ {1 / 2}} \sum_ {\sigma \in \mathfrak {S} _ {n}} x _ {\sigma (1)} \otimes \dots \otimes x _ {\sigma (n)} \quad (x _ {1},..., x _ {n} \text { in } \mathrm{E}).\tag{18}
\]

因此存在一个连续线性映射 \(\hat{\mathbf{S}}^{n}(u)\)（从 \(\hat{\mathbf{S}}^{n}(\mathrm{E})\) 到 \(\hat{\mathbf{S}}^{n}(\mathrm{F})\) 中）且仅存在一个，它使下面的图表交换

\[
\begin{array}{c c c} \hat {\mathbf {S}} ^ {n} (\mathrm{E}) & \xrightarrow {\phi_ {n , \mathrm{E}}} & \hat {\mathbf {T}} ^ {n} (\mathrm{E}) \\ \hat {\mathbf {S}} ^ {n} (u) \Big \downarrow & & \Big \downarrow \hat {\mathbf {T}} ^ {n} (u) \\ \hat {\mathbf {S}} ^ {n} (\mathrm{F}) & \xrightarrow {\phi_ {n , \mathrm{F}}} & \hat {\mathbf {T}} ^ {n} (\mathrm{F}) \end{array}
\]

我们现在证明公式

\[
\left\| \hat {\mathbf {S}} ^ {n} (u) \right\| = \| u \| ^ {n}.\tag{19}
\]

显然有 \(\| \hat{\mathbf{S}}^n (u)\| \leqslant \| \hat{\mathbf{T}}^n (u)\| = \| u\| ^n\)。另一方面，对一切 \(x\in \mathrm{E}\)，有 \(\hat{\mathbf{S}}^n (u)(x^n) = u(x)^n,\| x^n\| = (n!)^{1 / 2}\| x\| ^n\) 以及 \(\| u(x)^n\| = (n!)^{1 / 2}\| u(x)\| ^n\)，由此得

\[
\left\| \hat {\mathbf {S}} ^ {n} (u) \right\| \| x \| ^ {n} \geqslant \left\| u (x) \right\| ^ {n};
\]

由此立即得出 \(\| \hat{\mathbf{S}}^n (u)\| \geqslant \| u\| ^n\)，因而公式 (19) 成立。

显然我们有公式

\[
\hat {\mathbf {S}} ^ {n} (1 _ {\mathrm{E}}) = 1 _ {\hat {\mathbf {S}} ^ {n} (\mathrm{E})}\tag{20}
\]

\[
\hat {\mathbf {S}} ^ {n} (v \circ u) = \hat {\mathbf {S}} ^ {n} (v) \circ \hat {\mathbf {S}} ^ {n} (u) \quad \text { for } \quad v \in \mathscr {L} (\mathrm{F}; \mathrm{G}).\tag{21}
\]

最后，\(\hat{\mathbf{S}}^n (u)\) 在 \(\mathbf{S}^n (\mathrm{E})\) 上与 A, III, § 6, No.~2 中定义的线性映射 \(\mathbf{S}^n (u):\mathbf{S}^n (\mathrm{E})\to \mathbf{S}^n (\mathrm{F})\) 一致，因为它把 \(x_{1}\ldots x_{n}\) 变为 \(u(x_{1})\ldots u(x_{n})\)（对 E 中的每个 \(x_{1},\ldots ,x_{n}\)）。

例. --- * 1) 设 \(d \geqslant 1\) 为整数，\(\omega\) 为 \(\mathbf{R}^{d}\) 上关于 Lebesgue 测度 \(\mu\) 局部可积的正函数。设 E 为希尔伯特空间 \(\mathrm{L}^{2}(\mathbf{R}^{d}, \omega, \mu)\)，并令 \(\mathbf{S} = \mathbf{S}(\mathrm{E})\)。于是 \(\mathbf{S}\) 可与所有序列 \(f = (f_{n})_{n \geqslant 0}\) 组成的空间等同，其中每个 \(f_{n}\) 是 \((\mathbf{R}^{d})^{n}\) 上的函数，它关于 Lebesgue 测度 \(\mu \otimes \ldots \otimes \mu\)（\(n\) 个因子）可测，且在 \(n\) 个因子（\((\mathbf{R}^{d})^{n}\) 的）的置换下不变，并且满足

\[
\| \mathbf {f} \| ^ {2} = \sum_ {n = 0} ^ {\infty} n! \int_ {\mathbf {R} ^ {d}} \dots \int_ {\mathbf {R} ^ {d}} \left| f _ {n} (\mathbf {x} _ {1}, \dots , \mathbf {x} _ {n}) \right| ^ {2} \omega (\mathbf {x} _ {1}) \dots \omega (\mathbf {x} _ {n}) d \mathbf {x} _ {1} \dots d \mathbf {x} _ {n}\tag{22}
\]

为有限。范数 \(\| \mathbf{f}\|\)（\(\mathbf{S}\) 中的）由公式 (22) 定义。上面定义的希尔伯特空间 \(\mathbf{S}\) 称为对应于权 \(\omega_{*}\) 的对称 Fock 空间

* 2) 设 X 为分离的拓扑空间，\(\mu\) 为 X 上范数为 1 的正测度，E 为实希尔伯特空间 \(L_{\mathbb{R}}^{2}(X,\mu)\) 的希尔伯特子空间。若下列等价条件之一满足，则称 E 为高斯空间：

\begin{enumerate}
\def\labelenumi{\alph{enumi})}
\item
  对所有 \(f \in \mathrm{E}\)，有 \(\int_{\mathbf{x}} e^{if} d\mu = \exp(-\|f\|^2/2)\)；
\item
  对所有范数为 1 的 \(f \in E\)，测度 \(\mu\) 在 f 下的像是测度
\end{enumerate}

\[
(2 \pi) ^ {- 1 / 2} e ^ {- x ^ {2} / 2} d x.
\]

在 \(\mathbf{R}\) 上。

假定 E 是高斯空间。设 \(f_1, \ldots, f_n\) 为若干函数，其类 \(f_i\) 属于 E。我们用下式在 X 上定义一个函数 :\(f_1 \ldots f_n\):（称为 \(f_1, \ldots, f_n\) 的« Wick 乘积 »）

\[
: f _ {1} \dots f _ {n} := \sum_ {0 \leqslant 2 p \leqslant n} (- 1) ^ {p} \sum_ {\sigma \in I _ {p}} \prod_ {i = 1} ^ {p} \left\langle f _ {\sigma (2 i - 1)} \mid f _ {\sigma (2 i)} \right\rangle \prod_ {j = 2 p + 1} ^ {n} f _ {\sigma (j)},\tag{23}
\]

其中 \(\mathbf{I}_p\) 是由满足下列条件的置换 \(\sigma\)（集合 \(\{1,2,\dots,n\}\) 的）组成的集合

\[
\begin{array}{l} \sigma (1) <   \sigma (2), \dots , \sigma (2 p - 1) <   \sigma (2 p) \\ \sigma (1) <   \sigma (3) <   \dots <   \sigma (2 p - 1) \\ \sigma (2 p + 1) <   \sigma (2 p + 2) <   \dots <   \sigma (n). \end{array}
\]

于是存在一个同构 \(\phi\)（从 \(\hat{\mathbf{S}}(\mathrm{E})\) 到 \(\mathrm{L}_{\mathbb{R}}^{2}(\mathrm{X},\mu)\) 的一个希尔伯特子空间上），它把乘积 \(\dot{f}_1\dots \dot{f}_n\)（即 \(\dot{f}_1,\dots ,\dot{f}_n\) 在 \(\hat{\mathbf{S}}(\mathrm{E})\) 中算出的乘积）变为 \((:f_1\dots f_n:)\)。假定 X 是 Souslin 空间，且存在由类属于 E 且分离 X 的点的函数组成的可数族 \((f_{n})\)。这时 \(\phi\) 是从 \(\hat{\mathbf{S}}(\mathrm{E})\) 到 \(\mathrm{L}_{\mathbb{R}}^{2}(\mathrm{X},\mu)\) 上的同构。

\subsection{外希尔伯特幂}

设 E 为希尔伯特空间，n 为正整数。对每个置换 \(\sigma\in\mathfrak{S}_{n}\)，以 \(\varepsilon_{\sigma}\) 表示其符号差；令 \(a_{n}=\frac{1}{n!}\sum_{\sigma\in\mathfrak{S}_{n}}\varepsilon_{\sigma}p_{\sigma}\)（在 \(\mathcal{L}(\hat{\mathbf{T}}^{n}(\mathrm{E}))\) 中）（V, p.~29）。立即可见 \(a_{n}\) 是一个正交投影算子，其像 \(\overline{\mathbf{A}_{n}^{\prime}(\mathrm{E})}\) 是 \(\hat{\mathbf{T}}^{n}(\mathrm{E})\) 中空间 \(\mathbf{A}_{n}^{\prime}(\mathrm{E})\)（由所有 \(n(\mathrm{A},\mathrm{III},\S7,\mathrm{No}.4)\) 阶反对称张量组成的）的闭包。存在一个同构 \(\pi_{n}\)（从 \(\Lambda^{n}(\mathrm{E})\) 到 \(\mathbf{A}_{n}^{\prime}(\mathrm{E})\) 上），它由下式刻画

\[
\pi_ {n} (x _ {1} \wedge \ldots \wedge x _ {n}) = \mathbf {a} _ {n} (x _ {1} \otimes \ldots \otimes x _ {n}) = \frac {1}{n !} \sum_ {\sigma \in \mathfrak {S} _ {n}} \varepsilon_ {\sigma} x _ {\sigma (1)} \otimes \ldots \otimes x _ {\sigma (n)}\tag{24}
\]

其中 \(x_{1}, \ldots, x_{n}\) 属于 E。现在我们可以通过令下式成立而在 \(\Lambda^{n}(E)\) 上定义一个分离准希尔伯特空间结构

\[
\langle u | v \rangle = n! \left\langle \pi_ {n} (u) \mid \pi_ {n} (v) \right\rangle .\tag{25}
\]

更明确地，我们有（试与 A, III, § 11, No.~5 的公式 (30) 比较）。

\[
\langle x _ {1} \wedge \dots \wedge x _ {n} | y _ {1} \wedge \dots \wedge y _ {n} \rangle = \det (\langle x _ {i} | y _ {j} ^ {\prime} \rangle)\tag{26}
\]

其中 \(x_{1},\ldots ,x_{n}\) 与 \(y_{1},\ldots ,y_{n}\) 都属于 E。

设 \(\hat{\mathbf{A}}^n (\mathrm{E})\) 表示准希尔伯特空间 \(\mathbf{A}^n (\mathrm{E})\) 的完备化，\(\hat{\mathbf{A}} (\mathrm{E})\) 表示诸希尔伯特空间 \(\hat{\mathbf{A}}^n (\mathrm{E})\) 的外希尔伯特和。

例. --- * 采用 V, p.~32 例 1 的记号，我们可以把希尔伯特空间 \(\symbf{\Lambda}(\mathrm{E})\) 与所有可测函数序列 \((f_n)_{n\geqslant 0}\)（对它们，(22) 中定义的数 \(\| \mathbf{f}\|\) 为有限）组成的集合等同起来，其中每个函数 \(f_{n}\) 都是反对称的，即满足关系式

\[
f _ {n} (\mathbf {x} _ {\sigma (1)}, \dots , \mathbf {x} _ {\sigma (n)}) = \varepsilon_ {\sigma} f _ {n} (\mathbf {x} _ {1}, \dots , \mathbf {x} _ {n})
\]

其中 \(\sigma \in \mathfrak{S}_n\) 为任意置换。希尔伯特空间 \(\hat{\Lambda}(\mathrm{E})\) 称为对应于权 \(\omega\) 的反对称 Fock 空间。*

命题 5. --- 设 \((e_i)_{i \in I}\) 为希尔伯特空间 E 的标准正交基。给 I 赋以一个全序结构。于是所有元素 \(e_{i_1} \wedge \ldots \wedge e_{i_n}\)（对 \(i_1 < \cdots < i_n\)）组成的集合是 \(\hat{\Lambda}^n(E)\) 的标准正交基。

我们知道（A, III, § 7, No.~8），所说的这些元素构成向量空间 \(\Lambda^n (\mathrm{E}_0)\) 的基，其中 \(\mathrm{E}_0\) 是 E 中由向量 \(e_i\) 生成的向量子空间。再者，当 \(i_1 < \dots < i_n\) 时，内积 \(\langle e_{i_k}|e_{i_\ell}\rangle\) 组成的矩阵是 \(n\) 阶单位矩阵；由 (26)，有 \(\| e_{i_1}\wedge \ldots \wedge e_{i_n}\| = 1\)。最后，若 \((i_1,\dots,i_n)\) 与 \((j_1,\dots,j_n)\) 是 I 中元素的两个不同的严格递增序列，则存在一个元素 \(j_\ell\) 与 \(i_1,\ldots ,i_n\) 都不同，于是 \(\langle e_{i_k}|e_{j_\ell}\rangle = 0\)（对 \(1\leqslant k\leqslant n\)），从而由 (26) 得 \(\langle e_{i_1}\wedge \ldots \wedge e_{i_n}|e_{j_1}\wedge \ldots \wedge e_{j_n}\rangle = 0\)。换言之，元素 \(e_{i_1}\wedge \ldots \wedge e_{i_n}\)（对 \(i_1 < \dots < i_n\)）组成的族是标准正交的。

但 \(E_{0}\) 在 E 中稠密，且映射 \((x_{1},...,x_{n})\mapsto x_{1}\wedge\ldots\wedge x_{n}\)（从 \(E\times\cdots\times E\) 到 \(\Lambda^{n}(E)\) 中）是连续的。因此 \(\Lambda^{n}(E_{0})\) 在 \(\Lambda^{n}(E)\) 中稠密，由此即得命题 5。

推论. --- 假定希尔伯特空间 E 是两个正交子空间 M 与 N 的直和。从 \(\symbf{\Lambda}(\mathbf{M})\otimes\symbf{\Lambda}(\mathbf{N})\) 到 \(\symbf{\Lambda}(\mathbf{E})\) 上的典范同构 g（A, III, § 7, No.~7）以唯一的方式延拓为从 \(\hat{\symbf{\Lambda}}(\mathbf{M})\hat{\otimes}_{2}\hat{\symbf{\Lambda}}(\mathbf{N})\) 到 \(\hat{\symbf{\Lambda}}(\mathbf{E})\) 上的希尔伯特空间同构。

证明与命题 4 的推论（V, p.~31）的证明类似。

设 \(\mathbf{E}\) 与 \(\mathbf{F}\) 为两个希尔伯特空间，\(u \in \mathcal{L}(\mathbf{E};\mathbf{F})\)。与对称幂 \(\hat{\mathbf{S}}^n (\mathbf{E})\) 的情形（V, p.~32）一样，我们将证明线性映射 \(\symbf{\Lambda}^n (u)\)（从 \(\symbf{\Lambda}^n (\mathbf{E})\) 到 \(\symbf{\Lambda}^n (\mathbf{F})(\mathbf{A},\mathbf{III},\S 7,\mathbf{No}.4)\) 中）可延拓为连续线性映射 \(\hat{\symbf{\Lambda}}^n (u)\)（从 \(\hat{\symbf{\Lambda}}^n (\mathbf{E})\) 到 \(\hat{\symbf{\Lambda}}^n (\mathbf{F})\) 中）。我们有关系式

\[
\hat {\symbf {\Lambda}} ^ {n} (1 _ {\mathrm{E}}) = 1 _ {\hat {\symbf {\Lambda}} ^ {n} (\mathrm{E})},\tag{27}
\]

\[
\hat {\symbf {\Lambda}} ^ {n} (v \circ u) = \hat {\symbf {\Lambda}} ^ {n} (v) \circ \hat {\symbf {\Lambda}} ^ {n} (u) \quad \text { if } v \text { belongs   to } \mathscr {L} (\mathrm{F}; \mathrm{G}),\tag{28}
\]

\[
\left\| \hat {\symbf {\Lambda}} ^ {n} (u) \right\| \leqslant \| u \| ^ {n}.\tag{29}
\]

一般而言，公式 (29) 中的等号并不成立（TS, IV, § 6）。最后，我们有一个同构 \(\psi_{n} = \psi_{n,\mathrm{E}}\)（从 \(\hat{\mathbf{A}}^{n}(\mathrm{E})\) 到子空间 \(\overline{\mathbf{A}_{n}^{\prime}(\mathrm{E})}\)（\(\hat{\mathbf{T}}^{n}(\mathrm{E})\) 的）上），它由下式定义

\[
\psi_ {n} (x _ {1} \wedge \dots \wedge x _ {n}) = \frac {1}{(n !) ^ {1 / 2}} \sum_ {\sigma \in \mathfrak {S} _ {n}} \varepsilon_ {\sigma} x _ {\sigma (1)} \otimes \dots \otimes x _ {\sigma (n)}.\tag{30}
\]

\subsection{外乘法}

设 E 为希尔伯特空间。对每个整数 \(n \geqslant 0\)，以 \(\theta_{n}\) 表示从 \(\mathbf{T}^{n}(\mathbf{E})\) 到 \(\symbf{\Lambda}^{n}(\mathbf{E})\) 上的典范映射；于是

\[
\theta_ {n} (x _ {1} \otimes \ldots \otimes x _ {n}) = x _ {1} \wedge \ldots \wedge x _ {n}\tag{31}
\]

其中 \(x_{1}, \ldots, x_{n}\) 属于 E。设 \(p\) 与 \(q\) 为两个正整数；根据公式 (30) 与 (31)，我们有

\[
u \wedge v = \theta_ {p + q} \left(\frac {1}{(p !) ^ {1 / 2}} \psi_ {p} (u) \otimes \frac {1}{(q !) ^ {1 / 2}} \psi_ {q} (v)\right)\tag{32}
\]

其中 \(u \in \Lambda^p(\mathbf{E})\)，\(v \in \Lambda^q(\mathbf{E})\)。由于 \(\| \theta_n \| \leqslant (n!)^{1/2}\)，我们得到不等式

\[
\| u \wedge v \| \leqslant \left(\frac {(p + q) !}{p ! q !}\right) ^ {1 / 2} \| u \| \cdot \| v \|\tag{33}
\]

其中 \(u \in \symbf{\Lambda}^p(\mathrm{E})\)，\(v \in \symbf{\Lambda}^q(\mathrm{E})\)。因此，映射 \((u, v) \mapsto u \wedge v\) 可由连续性延拓为从 \(\hat{\symbf{\Lambda}}^p(\mathrm{E}) \times \hat{\symbf{\Lambda}}^q(\mathrm{E})\) 到 \(\hat{\symbf{\Lambda}}^{p + q}(\mathrm{E})\) 中的双线性映射，其范数至多等于 \(\left(\frac{(p + q)!}{p!q!}\right)^{1/2}\)（参看 V, p.~73, 练习 2）。我们仍把它记作 \((u, v) \mapsto u \wedge v\)。

命题 6. --- 设 E 为希尔伯特空间。我们有

\[
\| x \wedge u \| \leqslant \| x \| \cdot \| u \|\tag{34}
\]

其中 \(x\in \mathbf{E}\)，\(u\in \hat{\Lambda} (\mathrm{E})\)。

显然只需考虑 \(\| x\| = 1\) 的情形。

设 F 为 E 中由所有与 x 正交的向量组成的希尔伯特子空间。由于 E 是 F 与直线 K.x 的希尔伯特和，由 V, p.~34 的推论可知，映射 \((v,w)\mapsto v+x\wedge w\) 是从 \(\hat{\mathbf{A}}(\mathbf{F})\oplus\hat{\mathbf{A}}(\mathbf{F})\) 到 \(\hat{\mathbf{A}}(\mathbf{E})\) 上的希尔伯特空间同构。若 \(u=v+x\wedge w\)，其中 v, w 属于 \(\hat{\mathbf{A}}(\mathbf{F})\)，则有 \(x\wedge u=x\wedge v\)，从而 \(\|x\wedge u\|=\|v\|\leqslant(\|v\|^{2}+\|w\|^{2})^{1/2}=\|u\|\)。

推论 1. --- a) 设 \(x_1, \ldots, x_n\) 为希尔伯特空间 E 的元素。我们有

\[
\left\| x _ {1} \wedge \dots \wedge x _ {n} \right\| \leqslant \left\| x _ {1} \right\| \dots \left\| x _ {n} \right\|,\tag{35}
\]

其中等号仅当某个 \(x_{i}\) 为零或序列 \((x_{1},\dots,x_{n})\) 正交时才成立。

\begin{enumerate}
\def\labelenumi{\alph{enumi})}
\setcounter{enumi}{1}
\tightlist
\item
  设 \(x_{1}, \ldots, x_{n}, y_{1}, \ldots, y_{n}\) 为希尔伯特空间 E 的元素。我们有
\end{enumerate}

\[
\left| \det (\langle x _ {i}, y _ {j} \rangle) \right| \leqslant \| x _ {1} \| \dots \| x _ {n} \| \cdot \| y _ {1} \| \dots \| y _ {n} \|;\tag{36}
\]

若向量 \(x_{i}\) 与 \(y_{i}\) 都非零，则 (36) 中等号成立当且仅当 \((x_{1},\ldots ,x_{n})\) 与 \((y_{1},\dots,y_{n})\) 各自都是 \(\mathbf{E}\) 的同一向量子空间的正交基。

不等式 (35) 由命题 6 对 n 用归纳法得到；不等式 (36) 可由在 \(\hat{\Lambda}^{n}(E)\) 中应用 Cauchy-Schwarz 不等式以及 V, p.~33 的公式 (26) 推出。

假定序列 \((x_{1},\dots,x_{n})\) 正交。于是

\[
\left\| x _ {1} \wedge \dots \wedge x _ {n} \right\| ^ {2} = \det (\langle x _ {i} | x _ {j} \rangle) = \prod_ {i = 1} ^ {n} \left\| x _ {i} \right\| ^ {2}
\]

因为 \(\langle x_i|x_j\rangle = 0\)（对 \(i\neq j\)）。

现在假定向量 \(x_{1}, \ldots, x_{n}\) 都非零且不构成正交序列。由于 \(\| x_{1} \wedge \ldots \wedge x_{n} \|\) 对向量 \(x_{1}, \ldots, x_{n}\) 的依赖是对称的，我们可以假定 \(x_{1}\) 不与子空间 \(F\)（\(E\) 中由 \(x_{2}, \ldots, x_{n}\) 生成的）正交，且 \(F\) 不恰为 0。我们可以把 \(x_{1}\) 分解为 \(x_{1}' + y\) 的形式，其中 \(y \neq 0\) 属于 \(F\)，而 \(x_{1}'\) 与 \(F\) 正交；这时 \(\| x_{1}' \| < \| x_{1} \|\)。但是

\[
x _ {1} \wedge x _ {2} \wedge \ldots \wedge x _ {n} = x _ {1} ^ {\prime} \wedge x _ {2} \wedge \ldots \wedge x _ {n},
\]

因而

\[
\begin{array}{r l} \left\| x _ {1} \wedge \dots \wedge x _ {n} \right\| & \leqslant \left\| x _ {1} ^ {\prime} \right\| \left\| x _ {2} \right\| \dots \left\| x _ {n} \right\| \\ & <   \left\| x _ {1} \right\| \left\| x _ {2} \right\| \dots \left\| x _ {n} \right\|. \end{array}
\]

假定向量 \(x_{i}\) 与向量 \(y_{i}\) 都非零。关系式 (36) 中的等式等价于下列诸等式的合取

\[
\left| \left\langle x _ {1} \wedge \dots \wedge x _ {n} \mid y _ {1} \wedge \dots \wedge y _ {n} \right\rangle \right| = \| x _ {1} \wedge \dots \wedge x _ {n} \| \cdot \| y _ {1} \wedge \dots \wedge y _ {n} \|\tag{37}
\]

\[
\left\| x _ {1} \wedge \dots \wedge x _ {n} \right\| = \left\| x _ {1} \right\| \dots \left\| x _ {n} \right\|, \quad \left\| y _ {1} \wedge \dots \wedge y _ {n} \right\| = \left\| y _ {1} \right\| \dots \left\| y _ {n} \right\|.\tag{38}
\]

由证明的第一部分，等式 (38) 蕴含序列 \((x_{1},\ldots ,x_{n})\) 与 \((y_{1},\ldots ,y_{n})\) 各自都正交，这又蕴含 \(x_{1}\wedge \ldots \wedge x_{n}\neq 0\) 且 \(y_{1}\wedge \ldots \wedge y_{n}\neq 0\)。在这些条件下，关系式 (37) 蕴含存在一个标量 \(\lambda \neq 0\) 使 \(y_{1}\wedge \ldots \wedge y_{n} = \lambda x_{1}\wedge \ldots \wedge x_{n}\)（V, p.~3, 注 1）；换言之，\((x_{1},\dots,x_{n})\) 与 \((y_{1},\dots,y_{n})\) 是 E 的同一向量子空间的基（A, III, § 11, No.~13）。

推论 2. --- 设 \((a_{ij})_{1\leqslant i,j\leqslant n}\) 为一个复元素埃尔米特矩阵，其行列式为 D。假定不等式

\[
\sum_ {i, j = 1} ^ {n} a _ {i j} \overline {{z}} _ {i} z _ {j} \geqslant 0\tag{39}
\]

对所有复数 \(z_{1}, \ldots, z_{n}\) 成立。于是

\[
0 \leqslant \mathrm{D} \leqslant a _ {1 1} \dots a _ {n n}.\tag{40}
\]

假定 \(\mathbf{D}\) 非零；等式 \(\mathbf{D} = a_{11} \ldots a_{nn}\) 成立当且仅当 \(a_{ij} = 0\)（对所有 \(i \neq j\)）。

设 \(\Phi\) 为向量空间 \(\mathbf{C}^n\) 上由下式给出的埃尔米特型

\[
\Phi (\mathbf {z}, \mathbf {z} ^ {\prime}) = \sum_ {i, j = 1} ^ {n} a _ {i j} \overline {{{{z}}}} _ {i} z _ {j} ^ {\prime}
\]

其中 \(\mathbf{z} = (z_1, \dots, z_n)\) 与 \(\mathbf{z}' = (z_1', \dots, z_n')\) 属于 \(\mathbf{C}^n\)。根据假设，\(\Phi\) 是正的。

先假定 \(\Phi\) 是非退化的，即 D 非零。若 \((e_1, \ldots, e_n)\) 是 \(\mathbf{C}^n\) 的典范基，则 \(\Phi(\mathbf{e}_i, \mathbf{e}_j) = a_{ij}\)，于是推论 2 立即由推论 1 \(a)\)（令 \(x_i = \mathbf{e}_i\)）得出。

由于 \(a_{ii} = \Phi(e_i, e_i) \geqslant 0\)，当 \(D = 0\) 时不等式 (40) 也成立。

推论 3（« Hadamard 不等式 »）. --- 设 \((a_{ij})_{1 \leqslant i,j \leqslant n}\) 为一个复元素矩阵，其行列式为 D。令

\[
c _ {i} = \left(\sum_ {j = 1} ^ {n} | a _ {i j} | ^ {2}\right) ^ {1 / 2} \quad \text {for} \quad 1 \leqslant i \leqslant n,
\]

并令 \(m = \sup_{i,j}|a_{ij}|\)。于是我们有

\[
| \mathbf {D} | \leqslant c _ {1} \dots c _ {n} \leqslant m ^ {n}. n ^ {n / 2}.\tag{41}
\]

若 \(\mathbf{D} \neq 0\)，要使 \(|\mathbf{D}| = c_1 \ldots c_n\) 成立，必须且只需行 \(y_i = (a_{ij})_{1 \leqslant j \leqslant n}\)（矩阵 \((a_{ij})_{1 \leqslant i,j \leqslant n}\) 的行）是两两正交的向量。

给空间 \(C^{n}\) 赋以由下式定义的内积

\[
\langle \mathbf {z} | \mathbf {z} ^ {\prime} \rangle = \sum_ {i = 1} ^ {n} \overline {{z}} _ {i} z _ {i} ^ {\prime}.
\]

设 \((\mathbf{x}_{1},\ldots,\mathbf{x}_{n})\) 为 \(C^{n}\) 的典范基，\(y_{i}\) 为以 \(a_{ij}\)（\(1\leqslant j\leqslant n\)）为分量的向量。我们有 \(\|x_{i}\|=1\) 与 \(\|y_{i}\|=c_{i}\)（对 \(1\leqslant i\leqslant n\)）；此外 \(\langle x_{i}|y_{j}\rangle=a_{ji}\)。于是不等式 \(|D|\leqslant c_{1}\ldots c_{n}\) 及其取等条件都是 V, p.~35 推论 1 的特殊情形。显然有 \(c_{i}\leqslant m.n^{1/2}\)，因此 \(c_{1}\ldots c_{n}\leqslant m^{n}.n^{n/2}\)。

\section{希尔伯特空间中的若干类算子}

在本节中，\(I_{E}\) 表示希尔伯特空间 E 的恒等映射。两个线性映射的复合 \(v \circ u\) 通常记作 vu 或 v.u。

\subsection{伴随}

命题 1. --- 设 E 与 F 为两个希尔伯特空间。对每个映射 \(u \in \mathcal{L}(E; F)\)，存在唯一的映射 \(u^{*} \in \mathcal{L}(F; E)\)，使得

\[
\langle u (x) | y \rangle_ {\mathrm{F}} = \langle x | u ^ {*} (y) \rangle_ {\mathrm{E}}\tag{1}
\]

对所有 \(x \in E\) 与所有 \(y \in F\) 成立。映射 \(u \mapsto u^{*}\)（从 \(\mathcal{L}(\mathrm{E}, \mathrm{F})\) 到 \(\mathcal{L}(\mathrm{F}; \mathrm{E})\) 中）是双射、等距且半线性的（关于 K 的自同构 \(\xi \mapsto \bar{\xi}\)）。

设 \(\mathcal{S}(\mathrm{E},\mathrm{F})\) 为 \(\mathbf{E}\times \mathbf{F}\) 上所有连续半双线性型组成的空间，赋以范数

\[
\| \Phi \| = \sup _ {\| x \| \leqslant 1, \| y \| \leqslant 1} | \Phi (x, y) |.\tag{2}
\]

我们类似地定义空间 \(\mathcal{S}(\mathrm{F},\mathrm{E})\)。我们曾定义（V, p.~16, 推论 2）一个从 \(\mathcal{L}(\mathrm{E};\mathrm{F})\) 到 \(\mathcal{S}(\mathrm{F},\mathrm{E})\) 上的 Banach 空间同构，记作 \(u\mapsto\Phi_{u}\)，它由下式刻画

\[
\Phi_ {u} (y, x) = \langle y | u (x) \rangle_ {\mathrm{F}} (x \in \mathrm{E}, y \in \mathrm{F}).\tag{3}
\]

以类似的方式我们定义一个从 \(\mathcal{L}(\mathrm{F},\mathrm{E})\) 到 \(\mathcal{S}(\mathrm{E},\mathrm{F})\) 上的同构。最后，我们用下式定义映射 \(\Phi \mapsto \Phi^{*}\)（从 \(\mathcal{S}(\mathrm{F},\mathrm{E})\) 到 \(\mathcal{S}(\mathrm{E},\mathrm{F})\) 上）

\[
\Phi^ {*} (x, y) = \overline {{\Phi (y , x)}} \quad (x \in \mathrm{E}, y \in \mathrm{F}).\tag{4}
\]

这一映射是双射、半线性且等距的。但公式 (1) 可以改写为 \(\Phi_{u^{*}} = (\Phi_{u})^{*}\)，由此即得命题。

定义 1. --- 设 E 与 F 为两个希尔伯特空间。对每个连续线性映射 \(u: \mathbf{E} \to \mathbf{F}\)，由公式 (1) 定义的从 F 到 E 中的连续线性映射称为 \(u\) 的伴随，记作 \(u^*\)。

我们有

\[
(u + v) ^ {*} = u ^ {*} + v ^ {*}\tag{5}
\]

\[
(\lambda u) ^ {*} = \overline {{{{\lambda}}}} u ^ {*}\tag{6}
\]

\[
(u ^ {*}) ^ {*} = u\tag{7}
\]

\[
(1 _ {\mathrm{E}}) ^ {*} = 1 _ {\mathrm{E}}\tag{8}
\]

\[
(w u) ^ {*} = u ^ {*} w ^ {*};\tag{9}
\]

在所有这些公式中，u 与 v 属于 \(\mathcal{L}(\mathrm{E};\mathrm{F})\)，\(\lambda\) 属于 K，w 属于 \(\mathcal{L}(\mathrm{F};\mathrm{G})\)，其中 G 是一个希尔伯特空间。公式 (5) 与 (6) 表明 \(u\mapsto u^{*}\) 是半线性的。公式 (8) 是显然的。为证 (7)，我们对 (1) 的两边取共轭，得 \(\langle u^{*}(y)|x\rangle = \langle y|u(x)\rangle\)，这就证明 u 是 \(u^{*}\) 的伴随。最后，采用 (9) 的记号，对所有 \(z \in G\) 我们有

\[
\langle w (u (x)) | z \rangle = \langle u (x) | w ^ {*} (z) \rangle = \langle x | u ^ {*} (w ^ {*} (z)) \rangle ,
\]

因此 \(u^{*}w^{*}\) 是 wu 的伴随。

设 \(u: \mathrm{E} \to \mathrm{F}\) 为双射且连续的线性映射；于是它也是双连续的（I, p.~19, 推论 1）。由 (8) 与 (9) 我们立即推出 \(u^*\) 是双射且双连续的，并且

\[
(u ^ {- 1}) ^ {*} = (u ^ {*}) ^ {- 1}.\tag{10}
\]

命题 2. --- 对每个 \(u \in \mathcal{L}(\mathrm{E};\mathrm{F})\)，我们有

\[
\| u ^ {*} u \| = \| u u ^ {*} \| = \| u \| ^ {2} = \| u ^ {*} \| ^ {2}.\tag{11}
\]

由命题 1，\(\| u^{*}\| = \| u\|\)，因此 \(\| u^{*}u\| \leqslant \| u^{*}\| \cdot \| u\| \leqslant \| u\|^2\)。另一方面，

\[
\| u \| ^ {2} = \sup _ {\| x \| \leqslant 1} \| u (x) \| ^ {2} = \sup _ {\| x \| \leqslant 1} \langle u (x) | u (x) \rangle = \sup _ {\| x \| \leqslant 1} \langle x | u ^ {*} u (x) \rangle \leqslant \| u ^ {*} u \|,
\]

因此 \(\| u^{*}u\| = \| u\|^2\)。把 \(u\) 换成 \(u^{*}\)，得 \(\| uu^{*}\| = \| u^{*}\|^2\)，再由 \(\| u\| = \| u^{*}\|\) 即得 (11)。

设 \(E_{1}, \ldots, E_{n}\) 与 \(F_{1}, \ldots, F_{n}\) 为希尔伯特空间，且对 1 与 n 之间的每个整数 i，设 \(u_{i}\) 为从 \(E_{i}\) 到 \(F_{i}\) 中的连续线性映射。于是

\[
(u _ {1} \hat {\otimes} _ {2} \dots \hat {\otimes} _ {2} u _ {n}) ^ {*} = u _ {1} ^ {*} \hat {\otimes} _ {2} \dots \hat {\otimes} _ {2} u _ {n} ^ {*}.\tag{12}
\]

设 \(v\) 为连续线性映射 \(u_1 \hat{\otimes}_2 \ldots \hat{\otimes}_2 u_n\)，它从

\[
\mathbf {E} = \mathbf {E} _ {1} \hat {\otimes} _ {2} \dots \hat {\otimes} _ {2} \mathbf {E} _ {n} \quad \text {into} \quad \mathbf {F} = \mathbf {F} _ {1} \hat {\otimes} _ {2} \dots \hat {\otimes} _ {2} \mathbf {F} _ {n}
\]

而 w 为从 F 到 E 中的连续线性映射 \(u_{1}^{*}\hat{\otimes}_{2}\ldots\hat{\otimes}_{2}u_{n}^{*}\)。只需证明等式 \(\langle y|v(x)\rangle=\langle w(y)|x\rangle\)（对 \(x\in E\) 与 \(y\in F\)）。由线性性与连续性，可归结为 x 与 y 具有下列形式的情形

\[
x = x _ {1} \otimes \ldots \otimes x _ {n}, y = y _ {1} \otimes \ldots \otimes y _ {n}
\]

其中 \(x_{i} \in E_{i}\) 与 \(y_{i} \in F_{i}\)（对 \(1 \leqslant i \leqslant n\)）。于是由张量积中内积的定义（V, p.~27, 公式 (6)），得

\[
\langle y | v (x) \rangle = \prod_ {i = 1} ^ {n} \left\langle y _ {i} \mid u _ {i} \left(x _ {i}\right) \right\rangle = \prod_ {i = 1} ^ {n} \left\langle u _ {i} ^ {*} \left(y _ {i}\right) \mid x _ {i} \right\rangle = \left\langle w (y) \mid x \right\rangle .
\]

这就证明了我们的断言。

设 E 与 F 为两个希尔伯特空间，\(u \in \mathcal{L}(\mathrm{E}; \mathrm{F})\)，n 为正整数。若在公式 (12) 中令 \(u_{1} = \cdots = u_{n} = u\)，我们就得到如下结果：连续线性映射 \(\hat{\mathbf{T}}^{n}(u^{*})\)（从 \(\hat{\mathbf{T}}^{n}(\mathrm{F})\) 到 \(\hat{\mathbf{T}}^{n}(\mathrm{E})\) 中）是连续线性映射 \(\hat{\mathbf{T}}^{n}(u)\)（从 \(\hat{\mathbf{T}}^{n}(\mathrm{E})\) 到 \(\hat{\mathbf{T}}^{n}(\mathrm{F})\) 中）的伴随。公式

\[
\hat {\mathbf {S}} ^ {n} (u) ^ {*} = \hat {\mathbf {S}} ^ {n} (u ^ {*}), \quad \hat {\symbf {\Lambda}} ^ {n} (u) ^ {*} = \hat {\symbf {\Lambda}} ^ {n} (u ^ {*})\tag{13}
\]

可以仿照公式 (12) 来证明，只需用到 \(\hat{\mathbf{S}}^n (\mathrm{E})\)（V, p.~30, 公式 (15)）与 \(\hat{\Lambda}^n (\mathrm{E})\)（V, p.~33, 公式 (26)）中内积的定义。

注 1. --- 假定希尔伯特空间 E 不化为 0。我们把 \(\mathcal{L}(\mathbf{K};\mathbf{E})\) 与 E（借助映射 \(u\mapsto u(1)\)）等同；换言之，E 中的向量 \(x\) 与从 K 到 E 中的映射 \(\lambda\mapsto\lambda.x\) 等同。这时 \(x\) 的伴随是映射 \(x^{*}:E\to K\)（由 \(x^{*}(y)=\langle x|y\rangle\) 给出）。换句话说，\(x\mapsto x^{*}\) 是从 E 到其对偶上的典范半线性映射（V, p.~15）。

类似地，我们把数 \(\lambda \in \mathbf{K}\) 与 E 的自同态 \(\lambda. 1_{\mathrm{E}}\) 等同。这时 \(\lambda^{*}\) 恰好就是 \(\lambda\) 的共轭。

借助这些等同，我们可以定义乘积 \(t_{1} \ldots t_{n}\)，其中每个 \(t_{i}\) 或者是 K 中的数，或者是 E 中的向量，或者是属于 \(E'\) 的线性型，或者是 \(\mathcal{L}(E)\) 的元素，只要相邻的两个因子 \(t_{i}\) 与 \(t_{i+1}\) 从不属于下列类型之一：

\(\bullet\) \(xy\)，其中 \(x, y\) 都在 \(\mathbf{E}\) 中，或都在 \(\mathbf{E}'\) 中；

\(\bullet\) \(x\mathbf{A}\) 或 \(\mathbf{A}x'\)，其中 \(\mathbf{A} \in \mathcal{L}(\mathbf{E})\)，\(x \in \mathbf{E}\)，\(x' \in \mathbf{E}'\)。

我们有下列合成规则：

\begin{enumerate}
\def\labelenumi{\alph{enumi})}
\item
  结合性；
\item
  \(\mathbf{K}\) 的每个元素与其余所有因子交换；
\item
  我们有 \((t_1 \ldots t_n)^* = t_n^* \ldots t_1^*\)；换言之，乘积的伴随是诸伴随按相反顺序的乘积。此外 \(t^{**} = t\)。
\end{enumerate}

例如，设 \(x, y\) 属于 E，A 属于 \(\mathcal{L}(\mathrm{E})\)。于是 \(x^{*}y\) 表示内积 \(\langle x|y\rangle\)，而 \(x^{*}\mathrm{A}y\) 表示内积 \(\langle x|\mathrm{A}y\rangle\)。我们还有 \((\mathrm{A}^{*}x)^{*} = x^{*}\mathrm{A}^{**} = x^{*}\mathrm{A}\)，因此 \((\mathrm{A}^{*}x)^{*}y = x^{*}\mathrm{A}y\)，它可以解释为

\[
\langle \mathrm{A} ^ {*} x | y \rangle = \langle x | \mathrm{Ay} \rangle
\]

这与伴随的定义一致。我们注意到 \(yx^{*}\) 是 E 的自同态 \(z \mapsto y\langle x|z\rangle\)，因为由结合性，\(yx^{*}z\) 既可以解释为 \(y(x^{*}z)\)，也可以解释为 \(y.\langle x|z\rangle\)。

依照 Dirac \(^{1}\)，在数学物理的大多数著作中，E 的元素用符号 \(|x\rangle\) 表示，E' 的元素用 \(\langle t|\) 表示。内积写作 \(\langle x|y\rangle = \langle x|.|y\rangle\)，而乘积中的第一条禁用规则排除了记号 \(\rangle|\) 与 \(|\langle\) 的组合，例如 \(|x\rangle|y\rangle\)。

\(^{1}\) 参见 P. A. M. DIRAC, \emph{Quantum Mechanics}, Oxford University Press, New York, 1935.

命题 3. --- 设 E 与 F 为两个希尔伯特空间，\(u \in \mathcal{L}(\mathrm{E};\mathrm{F})\)。下列条件等价：

\begin{enumerate}
\def\labelenumi{(\roman{enumi})}
\item
  \(u\) 是拓扑向量空间同构，且其逆等于 \(u^{*}\)；
\item
  \(u\) 是满射且 \(u^{*}u = 1_{\mathrm{E}}\)；
\item
  \(u\) 是单射且 \(uu^{*} = 1_{\mathrm{F}}\)；
\item
  \(u\) 是赋范空间同构；
\item
  \(u\) 是希尔伯特空间同构。
\end{enumerate}

条件 (1) 意味着我们既有 \(u^{*}u = 1_{\mathrm{E}}\) 又有 \(uu^{*} = 1_{\mathrm{F}}\)。于是 (i)、(ii) 与 (iii) 的等价性由 E, II, § 3, No.~8, 命题 8 得出。(iv) 与 (v) 的等价性我们已经在（V, p.~5）中见过。最后，关系 \(u^{*}u = 1_{\mathrm{E}}\) 等价于 \(\langle x|u^{*}u(y)\rangle = \langle x|y\rangle\)，即 \(\langle u(x)|u(y)\rangle = \langle x|y\rangle\)（对 E 中所有 \(x,y\)），并且显然蕴含 \(u\) 是单射；这就证明了 (ii) 与 (v) 的等价性。

希尔伯特空间 E 的自同构也称为酉算子，即算子 \(u \in \mathcal{L}(\mathrm{E})\)（满足 \(uu^{*} = u^{*}u = 1_{E}\)）。

注 2. --- 关系 \(u^{*}u = 1_{\mathrm{E}}\) 并不刻画希尔伯特空间 E 的所有自同构。例如，设 \(\mathrm{E} = \ell^2 (\mathbf{N})\)，并设 \(u\) 由 \(u(x_{n}) = x_{n - 1}\)（\(n\geqslant 1\)）及 \(u(x)_0 = 0\) 定义。我们有 \(\| u(x)\| = \| x\|\)（对所有 \(x\in \mathrm{E}\)），即 \(u^{*}u = 1_{\mathrm{E}}\)，但 \(u\) 不是满射。

注 3. --- 伴随 \(u^{*}\) 的定义 (1) 也可以写成

\[
\langle y | u (x) \rangle = \langle u ^ {*} (y) | x \rangle ,
\]

或者，依 V, p.~15，写成

\[
\langle u (x), y ^ {*} \rangle = \langle x, (u ^ {*} (y)) ^ {*} \rangle .
\]

但我们也有 \(\langle u(x), y^* \rangle = \langle x, {}^t u(y^*) \rangle\)，因此我们可以用转置来表示伴随：

\[
(u ^ {*} (y)) ^ {*} = ^ {t} u (y ^ {*}).
\]

\subsection{部分等距线性映射}

定义 2. --- 设 E 与 F 为两个希尔伯特空间，\(u \in \mathcal{L}(E; F)\)。u 的核在 E 中的正交补称为 u 的初始子空间，u 的像在 F 中的闭包称为 u 的终子空间。从 E（相应地 F）到 u 的初始（相应地终）子空间上的正交投影算子称为 u 的初始（相应地终）正交投影算子。

设 \(\mathbf{P}\) 为 \(u\) 的初始子空间。由于 \(\mathbf{E}\) 是 \(\mathbf{P}\) 与 \(u\) 的核的直和，我们有 \(u(\mathbf{P}) = u(\mathbf{E})\)。

命题 4. --- (i) \(u^*\) 的初始（相应地终）子空间等于 \(u\) 的终（相应地初始）子空间。

\begin{enumerate}
\def\labelenumi{(\roman{enumi})}
\setcounter{enumi}{1}
\tightlist
\item
  假定 \(\mathbf{E} = \mathbf{F}\)。设 \(\mathbf{M}\) 为 \(\mathbf{E}\) 的闭向量子空间，\(\mathbf{M}^{\circ}\) 为其正交补。关系 \(u(\mathbf{M}) \subset \mathbf{M}\) 与 \(u^{*}(\mathbf{M}^{\circ}) \subset \mathbf{M}^{\circ}\) 等价。
\end{enumerate}

设 \(\mathrm{Q} = \overline{u(\mathrm{E})}\) 为 \(u\) 的终子空间。\(\mathrm{Q}^{\circ}\)（\(\mathrm{Q}\) 在 \(\mathrm{F}\) 中的正交补）由所有这样的向量 \(y\) 组成：\(\langle u(x)|y\rangle = 0\) 对一切 \(x\in \mathrm{E}\) 成立；这等价于：\(\langle x|u^{*}(y)\rangle = 0\) 对一切 \(x\in \mathrm{E}\) 成立，亦即 \(u^{*}(y) = 0\)。因此 \(\mathrm{Q}^{\circ} = \operatorname{Ker}u^{*}\)，而 \(\mathrm{Q}\) 是 \(u^{*}\) 的初始子空间。由于 \(u\) 是 \(u^{*}\) 的伴随，\(u^{*}\) 的终子空间也是 \(u\) 的初始子空间。这就证明了 (i)。

关系 \(u(\mathbf{M}) \subset \mathbf{M}\) 蕴含 \(u(\mathbf{M})\) 与 \(\mathbf{M}^{\circ}\) 正交，而关系 \(u^{*}(\mathbf{M}^{\circ}) \subset \mathbf{M}^{\circ}\) 蕴含 \(u^{*}(\mathbf{M}^{\circ})\) 与 \(\mathbf{M}\) 正交。但我们有 \(\langle u(x)|y\rangle = \overline{\langle u^{*}(y)|x\rangle}\)（对所有 \(x \in \mathbf{M}\) 与 \(y \in \mathbf{M}^{\circ}\)）；由此即得 (ii)。

我们注意到，利用 V, p.~41 的注 3，命题 4 可以由转置的一般性质（II, p.~51, 推论 2）推出。

定义 3. --- 设 E 与 F 为两个希尔伯特空间。称映射 \(u \in \mathcal{L}(\mathrm{E};\mathrm{F})\) 是部分等距的，如果对所有属于 u 的初始子空间的 x 有 \(\|u(x)\| = .\|x\|\)。

设 \(u \in \mathcal{L}(\mathrm{E};\mathrm{F})\)，\(\mathbf{N}\) 为其核，\(\mathbf{I}\) 为其像。说 \(u\) 是部分等距的，等于说线性映射 \(\tilde{u}: \mathrm{E/N} \to \mathrm{I}\)（由 \(u\) 导出的）是等距的（V, p.~13）。这时子空间 \(\mathbf{I}\)（\(\mathbf{F}\) 中的）是完备的，因而是闭的，并且是 \(u\) 的终子空间。于是，\(u\) 诱导出从 \(u\) 的初始子空间到其终子空间上的一个希尔伯特空间同构。

命题 5. --- 设 \(u \in \mathcal{L}(\mathrm{E};\mathrm{F})\)，\(\mathbf{P}\) 为其初始子空间，\(\mathbf{Q}\) 为其终子空间。以 \(p\)（相应地 \(q\)）表示 \(u\) 的初始（相应地终）正交投影算子。假定 \(u\) 是部分等距的。

\begin{enumerate}
\def\labelenumi{(\roman{enumi})}
\item
  映射 \(u^{*} \in \mathcal{L}(\mathbf{F};\mathbf{E})\) 是部分等距的，其初始子空间为 \(\mathbf{Q}\)，终子空间为 \(\mathbf{P}\)。这时从 \(\mathbf{P}\) 到 \(\mathbf{Q}\) 上的同构（由 \(u\) 诱导的）是从 \(\mathbf{Q}\) 到 \(\mathbf{P}\) 上的同构（由 \(u^{*}\) 诱导的）的逆。
\item
  我们有 \(u^{*}u = p\) 且 \(uu^{*} = q\)。
\end{enumerate}

根据命题 4 (i)，断言 (i) 是 (ii) 的推论。

我们现在证明 (ii)。由于 \(\mathbf{P}\) 包含 \(u^{*}\) 的像，映射 \(u^{*}u\) 把 \(\mathbf{E}\) 映到 \(\mathbf{P}\) 中。设 \(x \in \mathbf{E}\) 且 \(y \in \mathbf{P}\)，则

\[
\langle u ^ {*} u (x) | y \rangle = \langle u (x) | u (y) \rangle .
\]

若 x 属于 P，则由部分等距映射的定义有 \(\langle u(x)|u(y)\rangle = \langle x|y\rangle\)；若 x 属于 u 的核 N，则 \(u(x) = 0\)，从而 \(\langle u(x)|u(y)\rangle = 0\)，且由于 N 与 P 正交而有 \(\langle x|y\rangle = 0\)。因为 \(E = P \oplus N\)，所以在所有情形都有 \(\langle u^{*}u(x) - x|y\rangle = 0\)，于是 \(u^{*}u\) 是从 E 到 P 上的正交投影算子 p。\(uu^{*} = q\) 则由在上面的论证中把 u 与 \(u^{*}\) 互换得到。

命题 6. --- 对每个 \(u \in \mathcal{L}(\mathrm{E};\mathrm{F})\)，下列条件等价：

\begin{enumerate}
\def\labelenumi{(\roman{enumi})}
\item
  \(u\) 是部分等距的；
\item
  \(u^{*}\) 是部分等距的；
\item
  \(u^{*}u\) 是正交投影算子；
\item
  \(uu^{*}\) 是正交投影算子；
\item
  \(uu^{*}u = u;\)
\item
  \(u^{*}uu^{*} = u^{*}\)。
\end{enumerate}

由命题 5，(i) 与 (ii) 等价。

\begin{enumerate}
\def\labelenumi{(\roman{enumi})}
\item
  \(\Rightarrow\) (v)：假定 \(u\) 是部分等距的。于是由命题 5，\(u^{*}u\) 是 \(u\) 的初始正交投影算子。因此对每个 \(x \in \mathrm{E}\)，\(u^{*}u(x) - x\) 属于 \(u\) 的核，即 \(uu^{*}u(x) = u(x)\)。
\item
  \(\Rightarrow\) (iii)：假定 \(uu^{*}u = u\)，令 \(p = u^{*}u\)；则 \(p = p^{*}\) 且 \(p^2 = p\)。设 M（相应地 N）为 \(p\) 的像（相应地核）。对 \(x \in \mathbf{M}\) 与 \(y \in \mathbf{N}\)，我们有 \(\langle x|y\rangle = \langle p(x)|y\rangle = \langle x|p^{*}(y)\rangle = \langle x|p(y)\rangle = 0\)。由于 M 与 N 正交，\(p\) 是从 E 到 M 上的正交投影算子。
\item
  \(\Rightarrow\) (i)：假定 \(p = u^{*}u\) 是以 M 为像、以 N 为核的正交投影算子。
\end{enumerate}

对所有 \(x\in \mathbf{E}\)，我们有

\[
\left\| u (x) \right\| ^ {2} = \langle u ^ {*} u (x) | x \rangle = \langle p (x) | x \rangle .
\]

因此 \(u(x) = 0\)（对 \(x \in \mathbf{N}\)），且 \(\| u(x) \| = \| x \|\)（对 \(x \in \mathbf{M}\)），于是 \(u\) 是部分等距的，其核为 \(\mathbf{N}\)，初始子空间为 \(\mathbf{M}\)。

我们证明了 (i)、(iii) 与 (v) 的等价性。把 \(u\) 换成 \(u^*\)，即可推出 (ii)、(iv) 与 (vi) 的等价性。这就证明了命题 6。
